\label{chap:83-01}

% RESULTADO_RECUPERADO. Owner integro: sections/hmt/03_app_ortograma.tex
% Destino causal: capitulo 1; la jerarquia se subordina sin alterar el owner.
% Los dos bloques científicos del owner se activan en una única residencia.
\begingroup
\def\HMTAPPFormal{1}
\def\HMTAPPDevelopments{1}
\begin{HMTScientificOwnerSubordinated}
\input{sections/hmt/03_app_ortograma}
\end{HMTScientificOwnerSubordinated}
\endgroup

% Unidad U003. Redacción autónoma; tablas y censos reconstruidos desde 1063/live.

\section{Matriz aritmética de APP y estructura del anillo local}
\label{p12-83-c1-s1:chap:app-hojas-algebra-local}

La \cref{p12-83-c1-s2:chap:app-completacion-gnomonica} construyó un único soporte
\(X_9\). APP posee una única matriz aritmética principal: la multiplicación
nonádica publicada por raíz digital positiva. La suma no constituye una
segunda tabla APP; es la ley interna de acumulación que genera las filas y
columnas de la matriz multiplicativa y proporciona un observable auxiliar de
comparación. Esta distinción se mantendrá en toda la obra.

\subsection{Matriz multiplicativa nonádica y raíz digital}

Para \(a,b\in\mathcal D_9=\{1,\ldots,9\}\), la entrada de APP se obtiene
directamente por
\begin{equation}
 \boxed{\Pi(a,b)=\operatorname{dr}_9(ab)
 =1+((ab-1)\bmod9).}
 \label{p12-83-c1-s1:eq:app-matriz-multiplicativa-directa}
\end{equation}
La única matriz APP es, por tanto,
\[
\begin{array}{c|ccccccccc}
\Pi&1&2&3&4&5&6&7&8&9\\\hline
1&1&2&3&4&5&6&7&8&9\\
2&2&4&6&8&1&3&5&7&9\\
3&3&6&9&3&6&9&3&6&9\\
4&4&8&3&7&2&6&1&5&9\\
5&5&1&6&2&7&3&8&4&9\\
6&6&3&9&6&3&9&6&3&9\\
7&7&5&3&1&8&6&4&2&9\\
8&8&7&6&5&4&3&2&1&9\\
9&9&9&9&9&9&9&9&9&9
\end{array}
\]
Cada fila puede reconstruirse también por adición reiterada de su índice,
pero esa recurrencia describe la ley interna de la matriz y no una segunda
hoja. Las filas y columnas de índices \(1,2,4,5,7,8\) recorren las nueve
clases; las de índices \(3,6,9\) son exactamente las que contraen el soporte.

En la carta residual, el conjunto de unidades es
\[
 R_9^\times=(\mathbb Z/9\mathbb Z)^\times
   =\{[1],[2],[4],[5],[7],[8]\}.
\]
El ideal maximal es
\[
 \mathfrak m=(3)=\{[0],[3],[6]\}.
\]

\begin{theorem}[Clasificación de filas multiplicativas]
\label{p12-83-c1-s1:thm:app-clasificacion-filas-producto}
Para \(a\in\mathcal D_9\), la aplicación
\[
 m_a:\mathcal D_9\longrightarrow\mathcal D_9,
 \qquad b\longmapsto\Pi(a,b)
\]
es biyectiva si y sólo si \([a]_9\in R_9^\times\). Si
\([a]_9\in\mathfrak m\setminus\{0\}\), su imagen tiene tres elementos. Si
\([a]_9=[0]_9\), su imagen tiene un solo elemento.
\end{theorem}

\begin{proof}
La reducción residual de \(m_a\) es la multiplicación por \([a]_9\). En un
anillo finito, esa multiplicación es biyectiva exactamente cuando \([a]_9\)
es una unidad. Si \([a]_9=[3]\) o \([6]\), su imagen es el ideal
\(\mathfrak m\), de cardinal tres. Si \([a]_9=[0]\), toda entrada se envía
a la clase nula, publicada como \(9\).
\end{proof}

\begin{proposition}[Latinidad de la ley aditiva interna]
\label{p12-83-c1-s1:prop:app-latinidad-aditiva}
El observable auxiliar
\[
 \Sigma(a,b)=\operatorname{dr}_9(a+b)
\]
actúa por traslaciones: para cada \(a,c\in\mathcal D_9\) existe un único
\(b\in\mathcal D_9\) tal que \(\Sigma(a,b)=c\), y análogamente en la otra
coordenada.
\end{proposition}

\begin{proof}
En la carta residual, \([a+b]_9=[c]_9\) tiene la solución única
\([b]_9=[c-a]_9\). La sección positiva transporta esa clase a un único
elemento de \(\mathcal D_9\).
\end{proof}

La diferencia de geometría de fibras queda así tipada sin duplicar la matriz
APP: la ley aditiva interna sólo traslada, mientras la matriz multiplicativa
permuta sobre las unidades y contrae sobre el radical.

\begin{proposition}[Estructura lineal del radical]
\label{p12-83-c1-s1:prop:app-radical-lineal}
La aplicación
\[
 M_3:R_9\longrightarrow R_9,
 \qquad x\longmapsto3x
\]
satisface
\[
 \ker M_3=\operatorname{im}M_3=\mathfrak m.
\]
El mapa inducido
\[
 R_9/\mathfrak m\longrightarrow\mathfrak m,
 \qquad [x]\longmapsto3x
\]
es un isomorfismo de espacios vectoriales sobre \(\mathbb F_3\), pero no un
isomorfismo de anillos.
\end{proposition}

\begin{proof}
La congruencia \(3x\equiv0\pmod9\) equivale a \(x\equiv0\pmod3\), de
donde el núcleo \(\mathfrak m\). La imagen es
\(\{[0],[3],[6]\}=\mathfrak m\). El primer teorema de isomorfía da el
isomorfismo lineal. No es un isomorfismo de anillos porque
\(\mathfrak m^2=0\) en \(R_9\), mientras el cociente
\(R_9/\mathfrak m\cong\mathbb F_3\) tiene unidad no nula.
\end{proof}

La sección positiva publica el ciclo radical no trivial como
\[
 369\longmapsto693\longmapsto936\longmapsto369,
\]
según la posición desde la que se lea la terna. Esta notación conserva el
orden de las coordenadas; no convierte el ideal en tres números enteros
independientes.

\subsection{Levantamientos y cociclo de acarreo}

Sea \(s_0:R_9\to\{0,\ldots,8\}\) la sección residual normalizada. El
acarreo binario de la suma es
\[
 \kappa_0(a,b)
 =\frac{s_0(a)+s_0(b)-s_0(a+b)}9.
\]

\begin{lemma}[Identidad cocíclica]
\label{p12-83-c1-s1:lem:app-identidad-cociclo-acarreo}
Para \(a,b,c\in R_9\),
\[
 \kappa_0(a,b)+\kappa_0(a+b,c)
 =\kappa_0(b,c)+\kappa_0(a,b+c).
\]
\end{lemma}

\begin{proof}
Multiplicar por nueve y desarrollar cualquiera de los dos miembros produce
\[
 s_0(a)+s_0(b)+s_0(c)-s_0(a+b+c).
\]
Por tanto, ambos son iguales.
\end{proof}

Definimos \(\chi_0(a)=1\) si \(a=[0]_9\) y \(0\) en otro caso. Como
\[
 s_+(a)=s_0(a)+9\chi_0(a),
\]
los acarreos de las dos secciones se relacionan mediante
\[
 \kappa_+(a,b)=
 \kappa_0(a,b)+\chi_0(a)+\chi_0(b)-\chi_0(a+b).
\]
El cambio de representante modifica el cociclo por una coborde. Ésta es la
forma exacta de transportar la memoria cuando la clase nula se publica como
\(9\) en vez de \(0\).

\subsection{Censos integrados de residuos y cocientes}

Sumando sobre todos los pares \((i,j)\in\mathcal D_9^2\), se obtiene
\[
 \sum_{i,j}(i+j)=810,
 \qquad
 \sum_{i,j}ij=2025.
\]
Los residuos positivos satisfacen
\[
 \sum_{i,j}R^+_{i,j}=405,
 \qquad
 \sum_{i,j}R^\times_{i,j}=459.
\]
Por las identidades levantadas,
\[
 \sum_{i,j}Q^+_{i,j}=\frac{810-405}{9}=45,
 \qquad
 \sum_{i,j}Q^\times_{i,j}=\frac{2025-459}{9}=174.
\]

\begin{theorem}[Separación exacta del exceso suma--producto]
\label{p12-83-c1-s1:thm:app-separacion-exceso}
La diferencia integrada entre producto y suma se descompone de manera única
en una parte publicada y una parte de cociente:
\[
 \boxed{
  2025-810
  =(459-405)+9(174-45)
  =54+9\cdot129.}
\]
\end{theorem}

\begin{proof}
Restamos, término a término, las identidades
\(ij=R^\times_{i,j}+9Q^\times_{i,j}\) e
\(i+j=R^+_{i,j}+9Q^+_{i,j}\), y sumamos sobre los ochenta y un pares. Los
censos anteriores dan la igualdad. La unicidad procede de la división
euclídea en cada par.
\end{proof}

El valor \(54\) no es la diferencia completa entre dos supuestas matrices.
Es el defecto visible entre la matriz multiplicativa y el acumulador aditivo
en la carta positiva. El término \(9\cdot129\) es la contribución
que permanece en los cocientes. Omitirlo confundiría publicación con
genealogía.

\subsection{Bloque central de \(2\) posiciones}

Para \(k\in\mathbb Z/9\mathbb Z\), consideramos el bloque simétrico
\[
 B_k=
 \begin{pmatrix}
  k^2&k(k+1)\\
  k(k+1)&(k+1)^2
 \end{pmatrix}\pmod9.
\]
Las diagonales coinciden exactamente cuando
\[
 k^2\equiv(k+1)^2\pmod9
 \quad\Longleftrightarrow\quad
 2k+1\equiv0\pmod9.
\]
Como \(2^{-1}=5\pmod9\), la única solución es \(k=4\). Así aparece el
bloque
\[
 B_c=
 \begin{pmatrix}7&2\\2&7\end{pmatrix}.
\]
Sea
\[
 R=\begin{pmatrix}0&1\\1&0\end{pmatrix},
 \qquad
 P_\pm=\frac12(I\pm R).
\]
Entonces
\[
 B_c=7I+2R=9P_++5P_-.
\]

\begin{theorem}[Unicidad y espectro del bloque central]
\label{p12-83-c1-s1:thm:app-bloque-central}
El bloque \(B_c\) es el único bloque adyacente de la familia \(B_k\) con
diagonales iguales. Sus autovalores en los subespacios simétrico y
antisimétrico son, respectivamente, \(9\) y \(5\).
\end{theorem}

\begin{proof}
La congruencia anterior prueba la unicidad. Puesto que \(R\) actúa por
\(+1\) sobre \(\operatorname{im}P_+\) y por \(-1\) sobre
\(\operatorname{im}P_-\), la matriz \(7I+2R\) actúa por \(9\) y \(5\),
respectivamente.
\end{proof}

Las dos lecturas ordenadas de sus cifras dan
\[
 72+27=99,
 \qquad
 72-27=45=\det B_c,
\]
y las dos diagonales dan
\[
 77+22=99,
 \qquad
 77-22=55=54+1.
\]
Estas identidades son controles internos del bloque; no definen por sí solas
el defecto global.

\subsection{Curvatura mixta de las \(2\) leyes}

En la carta residual escribimos
\[
 S(x,y)=x+y,
 \qquad
 P(x,y)=xy.
\]
Para una función \(T:X_9\to R_9\), sean
\[
 \Delta_xT(x,y)=T(x+1,y)-T(x,y),
 \qquad
 \Delta_yT(x,y)=T(x,y+1)-T(x,y).
\]
Definimos la curvatura ordenada de un par de leyes por
\[
 \mathcal F(T_x,T_y)
 :=\Delta_x\Delta_yT_y-\Delta_y\Delta_xT_x.
\]

\begin{proposition}[Orientaciones de la composición mixta]
\label{p12-83-c1-s1:prop:app-curvatura-mixta}
Se cumplen
\[
 \mathcal F(S,S)=0,
 \qquad
 \mathcal F(P,P)=0,
\]
y, para las composiciones mixtas,
\[
 \mathcal F(S,P)=+1,
 \qquad
 \mathcal F(P,S)=-1
 \quad\text{en }R_9.
\]
\end{proposition}

\begin{proof}
Las segundas diferencias mixtas de la ley lineal \(S\) se anulan. Para
\(P(x,y)=xy\), ambas diferencias mixtas puras valen uno y se cancelan al
compararse con el mismo orden. Al combinar una ley lineal y una bilineal,
queda una única unidad cuyo signo cambia al invertir el orden de las hojas.
\end{proof}

La curvatura no crea otra matriz APP. Registra que el orden
\(\Sigma\to\Pi\) y el orden \(\Pi\to\Sigma\) son orientaciones opuestas,
mientras los sectores puros permanecen planos. Esta tríada
\(-1,0,+1\) será tipada en la unidad TRIT.

\subsection{Obstrucciones de la matriz aritmética de APP: traslaciones, radical local, bloque central, curvatura mixta y conservación del acarreo}

\begin{criterion}[Criterios de refutación de la matriz APP]
La unidad queda refutada si:
\begin{enumerate}
 \item la ley aditiva interna deja de actuar por traslaciones;
 \item una fila multiplicativa no unitaria se declara biyectiva;
 \item se identifica \(R_9/\mathfrak m\) con \(\mathfrak m\) como anillo;
 \item el exceso \(1215\) se sustituye por \(54\) sin conservar cocientes;
 \item el bloque central se escoge sin resolver \(2k+1=0\pmod9\);
 \item la inversión del orden mixto no cambia el signo de la curvatura;
 \item el acarreo se elimina al cambiar entre las secciones \(s_0\) y
       \(s_+\).
\end{enumerate}
\end{criterion}

APP está ahora construida como soporte, categoría de trayectorias, matriz
multiplicativa nonádica, ley aditiva interna, radical local, cociclo de
acarreo y orientación mixta. La unidad
siguiente no resumirá estas estructuras: tomará la tríada de orientaciones
que acaba de emerger y construirá el estado completo del TRIT con sus tres regímenes
cuadráticos.


% Unidad U002. Redacción autónoma a partir del generador integral 1063/live.

\section{Completación gnomónica y soporte aritmético de APP}
\label{p12-83-c1-s2:chap:app-completacion-gnomonica}

La presente unidad construye el soporte de APP a partir de la célula
nonádica de la \cref{p12-83-c1-s4:chap:segmentos-marcas-intersticios}. No se presupone un
plano continuo, una retícula infinita ni una constante numérica. El objeto de
partida es finito: dos copias ordenadas de la carta positiva de
\(\mathbb Z/9\mathbb Z\). La completación de sus intersticios interiores
produce exactamente ochenta y un estados, y la ley de traslación residual
convierte ese soporte cuadrado en un toro discreto.

En toda la unidad, APP designa una estructura matemática y no una imagen.
La sigla APP nombra en lo sucesivo esta estructura; las propiedades que
siguen dependen de sus objetos y morfismos, no de una expansión verbal de la
sigla.

\subsection{Correspondencia entre \(2\) cartas de una misma clase residual}

Definimos
\[
 \mathcal D_9:=\{1,2,\ldots,9\},
 \qquad
 \mathcal D_8:=\{1,2,\ldots,8\}.
\]
La aplicación de publicación residual es
\[
 \lambda_9:\mathcal D_9\longrightarrow\mathbb Z/9\mathbb Z,
 \qquad
 \lambda_9(a)=[a]_9.
\]
Por tanto, \(\lambda_9(9)=[0]_9\). Su inversa como aplicación de
conjuntos es la sección positiva
\[
 s_+([a]_9)=
 \begin{cases}
  a,&[a]_9\neq[0]_9,\quad 1\leq a\leq8,\\
  9,&[a]_9=[0]_9.
 \end{cases}
\]
Las dos expresiones \(9\) y \([0]_9\) no son dos fases: son dos
representantes del mismo elemento en cartas diferentes.

Para dos coordenadas escribiremos
\[
 \lambda_9^{\times2}:\mathcal D_9^2
       \longrightarrow X_9:=(\mathbb Z/9\mathbb Z)^2,
 \qquad
 (a,b)\longmapsto([a]_9,[b]_9).
\]
Esta aplicación es una biyección de conjuntos. La carta
\(\mathcal D_9^2\) hace visibles las marcas positivas; el grupo \(X_9\)
hace visibles la suma residual y las traslaciones.

\begin{lemma}[Cambio de carta sin pérdida de estado]
\label{p12-83-c1-s2:lem:app-cambio-carta-sin-perdida}
Las aplicaciones \(\lambda_9^{\times2}\) y
\(s_+^{\times2}\) son inversas. En particular, la escritura positiva y la
escritura residual contienen exactamente los mismos ochenta y un estados.
\end{lemma}

\begin{proof}
En cada coordenada se cumplen
\(\lambda_9\circ s_+=\operatorname{id}_{\mathbb Z/9\mathbb Z}\) y
\(s_+\circ\lambda_9=\operatorname{id}_{\mathcal D_9}\). El producto
cartesiano de estas identidades proporciona las dos identidades requeridas.
\end{proof}

\subsection{Núcleo interior y complemento gnomónico}

El núcleo interior es
\[
 \mathcal N_8:=\mathcal D_8\times\mathcal D_8,
 \qquad |\mathcal N_8|=8^2=64.
\]
En el plano de índices, \(\mathcal N_8\) es el cuadrado que no utiliza la
marca de cierre en ninguna coordenada. Para publicar también la clase nula
se debe completar la novena fila y la novena columna. Definimos el
complemento gnomónico disjunto por
\[
 \mathcal G_9:=
  (\{9\}\times\mathcal D_9)
  \sqcup
  (\mathcal D_8\times\{9\}).
\]
La esquina \((9,9)\) pertenece a la primera pieza; por eso la segunda usa
\(\mathcal D_8\) y no \(\mathcal D_9\). En consecuencia,
\[
 |\mathcal G_9|=9+8=17,
 \qquad
 \mathcal D_9^2=\mathcal N_8\sqcup\mathcal G_9,
 \qquad
 81=64+17.
\]

\begin{proposition}[Minimalidad de la completación]
\label{p12-83-c1-s2:prop:app-minimalidad-completacion}
Sea \(Y\) un subconjunto de \(\mathcal D_9^2\) que contiene
\(\mathcal N_8\) y cuya proyección sobre cada coordenada es todo
\(\mathcal D_9\). Si, además, \(Y\) contiene todos los pares necesarios
para trasladar una coordenada mientras la otra permanece arbitraria,
entonces \(\mathcal G_9\subseteq Y\). Por tanto, la adición de diecisiete
estados es la completación mínima compatible con las dos traslaciones.
\end{proposition}

\begin{proof}
Para trasladar la primera coordenada a través de la clase de cierre mientras
la segunda permanece en un valor arbitrario \(b\in\mathcal D_9\), son
necesarios todos los pares \((9,b)\). Se obtienen nueve estados. Para
trasladar la segunda coordenada a través del cierre mientras la primera toma
un valor arbitrario \(a\in\mathcal D_8\), son necesarios los ocho pares
\((a,9)\) restantes. La esquina \((9,9)\) ya fue contada en la primera
familia. En total son necesarios \(9+8=17\) estados, exactamente los de
\(\mathcal G_9\).
\end{proof}

La hipótesis de compatibilidad con las dos traslaciones es esencial. Bastaría
añadir una sola marca por coordenada para que ambas proyecciones fueran
sobreyectivas, pero ese conjunto no contendría los estados fronterizos
necesarios para mover una coordenada con la otra fija. APP requiere el
producto completo, no dos ejes desconectados.

\subsection{Realización exacta del intervalo intersticial mediante celdas APP y coordenadas de frontera}

Sea \(P_9\) el camino ordenado con nueve marcas y ocho aristas interiores
\[
 E(P_9)=\{e_1,\ldots,e_8\}.
\]
El cierre de los extremos incorpora una arista \(e_9\) y produce el ciclo
\(C_9\):
\[
 E(C_9)=E(P_9)\sqcup\{e_9\}.
\]
Los pares de intersticios interiores forman \(E(P_9)^2\). Los pares de
intersticios del ciclo completo forman \(E(C_9)^2\). La diferencia es
\begin{align}
 E(C_9)^2\setminus E(P_9)^2
   &=(\{e_9\}\times E(C_9))
     \sqcup(E(P_9)\times\{e_9\}).
\label{p12-83-c1-s2:eq:app-complemento-intersticial}
\end{align}
La descomposición es disjunta porque la esquina \((e_9,e_9)\) se reserva a
la primera pieza.

\begin{theorem}[Equivalencia del modelo de marcas y el modelo de
intersticios]
\label{p12-83-c1-s2:thm:app-equivalencia-modelos}
La aplicación
\[
 \beta:E(C_9)^2\longrightarrow\mathcal D_9^2,
 \qquad
 \beta(e_i,e_j)=(i,j),
\]
es biyectiva. Sus restricciones satisfacen
\[
 \beta(E(P_9)^2)=\mathcal N_8,
 \qquad
 \beta(E(C_9)^2\setminus E(P_9)^2)=\mathcal G_9.
\]
Por tanto, la identidad \(81=64+17\) cuenta simultáneamente pares de
marcas publicadas y pares de intersticios cerrados.
\end{theorem}

\begin{proof}
Cada arista de \(C_9\) tiene un índice único en \(\mathcal D_9\). La
aplicación \((i,j)\mapsto(e_i,e_j)\) es la inversa explícita de \(\beta\).
Si \(i,j\leq8\), el par pertenece a \(E(P_9)^2\) y su imagen pertenece a
\(\mathcal D_8^2\). Si al menos uno de los índices es nueve, el par pertenece
al complemento de \cref{p12-83-c1-s2:eq:app-complemento-intersticial}; separar el caso
\(i=9\) del caso \(i\leq8,j=9\) produce exactamente las dos piezas de
\(\mathcal G_9\).
\end{proof}

\subsection{Cierre toroidal y traslaciones}

Sean
\[
 \mathbf e_1=([1]_9,[0]_9),
 \qquad
 \mathbf e_2=([0]_9,[1]_9).
\]
El conjunto de direcciones elementales es
\[
 \mathscr D_{\rm APP}:=
 \{+\mathbf e_1,-\mathbf e_1,+\mathbf e_2,-\mathbf e_2\}.
\]
Definimos el grafo de Cayley orientado
\[
 \Gamma_{\rm APP}:=\operatorname{Cay}(X_9,\mathscr D_{\rm APP}).
\]
Sus vértices son los ochenta y un elementos de \(X_9\). Para cada
\(x\in X_9\) y \(d\in\mathscr D_{\rm APP}\) existe una arista orientada
\(x\xrightarrow{d}x+d\). La arista inversa lleva la etiqueta \(-d\).

En la carta positiva, la traslación unitaria se escribe mediante el sucesor
publicado y su acarreo:
\[
 u_+(a)=
 \begin{cases}
  a+1,&a<9,\\
  1,&a=9,
 \end{cases}
 \qquad
 \kappa_+(a)=
 \begin{cases}
  0,&a<9,\\
  1,&a=9.
 \end{cases}
\]
La identidad levantada es
\begin{equation}
 a+1=u_+(a)+9\kappa_+(a).
\label{p12-83-c1-s2:eq:app-sucesor-publicado-acarreo}
\end{equation}
Al aplicar \(\lambda_9\), el término de acarreo se anula y queda
\[
 \lambda_9(u_+(a))=\lambda_9(a)+[1]_9.
\]

\begin{lemma}[Compatibilidad entre cierre, residuo y elevación]
\label{p12-83-c1-s2:lem:app-compatibilidad-cierre-residuo}
La adición de \(e_9\), la identificación \(9\sim[0]_9\) y la traslación
\(x\mapsto x+\mathbf e_i\) describen tres operaciones compatibles en
dominios distintos. El cociente residual publica el retorno; el acarreo de
\cref{p12-83-c1-s2:eq:app-sucesor-publicado-acarreo} conserva la elevación que el
cociente omite.
\end{lemma}

\begin{proof}
El cierre topológico incorpora la arista que une los extremos del camino. La
carta residual identifica el extremo publicado con la clase nula. Si
\(a<9\), el sucesor tiene acarreo cero; si \(a=9\), la igualdad levantada
es \(10=1+9\). Ambos casos descienden a la misma suma residual, pero el
valor de \(\kappa_+\) los distingue.
\end{proof}

\begin{corollary}[Soporte toroidal]
\label{p12-83-c1-s2:cor:app-soporte-toroidal}
Identificar lados opuestos de la carta \(\mathcal D_9^2\) produce una
realización geométrica del grupo \(X_9\). Cada traslación elemental está
definida en todo vértice y posee inversa.
\end{corollary}

\subsection{Categoría de trayectorias APP: objetos, morfismos de prolongación y composición compatible}

\begin{definition}[Trayectoria orientada]
Una trayectoria de longitud \(n\geq0\) es un par
\[
 \gamma=(x_0;d_1,\ldots,d_n),
 \qquad
 x_0\in X_9,\qquad d_k\in\mathscr D_{\rm APP}.
\]
Sus vértices son
\[
 x_k=x_0+\sum_{r=1}^k d_r,
 \qquad 0\leq k\leq n.
\]
La fuente es \(s(\gamma)=x_0\) y el blanco es
\(t(\gamma)=x_n\). Para \(n=0\) se obtiene la trayectoria identidad en
\(x_0\).
\end{definition}

Si \(t(\gamma)=s(\delta)\), su concatenación es
\[
 \delta\circ\gamma
  =(x_0;d_1,\ldots,d_n,e_1,\ldots,e_m),
\]
donde \(e_1,\ldots,e_m\) son las direcciones de \(\delta\). La
concatenación es asociativa y las trayectorias de longitud cero son
identidades. Obtenemos así la categoría libre
\[
 \mathsf{Path}(\Gamma_{\rm APP}).
\]

\begin{proposition}[Finitud del soporte y numerabilidad de las trayectorias]
\label{p12-83-c1-s2:prop:app-soporte-finito-trayectorias-numerables}
El conjunto de estados de APP es finito, de cardinal \(81\). Para una fuente
fijada existen \(4^n\) palabras de direcciones de longitud \(n\). La unión
de todas las trayectorias finitas es numerable.
\end{proposition}

\begin{proof}
La primera afirmación procede de \(|X_9|=9^2\). En cada uno de los \(n\)
pasos hay cuatro elecciones independientes, de donde \(4^n\). Finalmente,
\(\bigcup_{n\geq0}X_9\times\mathscr D_{\rm APP}^n\) es una unión numerable
de conjuntos finitos.
\end{proof}

Esta proposición separa dos niveles que no se identificarán: APP posee un
soporte finito, pero admite trayectorias de cualquier longitud. La
prolongación de una ruta no crea nuevos vértices del soporte; crea nueva
genealogía sobre ellos.

\subsection{Matriz multiplicativa y ley aditiva interna sobre un único soporte}

Para \(n\in\mathbb Z_{>0}\), la raíz digital positiva es
\[
 \operatorname{dr}_9(n):=1+((n-1)\bmod9).
\]
Sobre la carta positiva definimos la matriz principal y la ley interna
\[
 \Sigma(a,b):=\operatorname{dr}_9(a+b),
 \qquad
 \Pi(a,b):=\operatorname{dr}_9(ab).
\]
Ambas aplicaciones tienen el mismo dominio \(\mathcal D_9^2\) y el mismo
codominio \(\mathcal D_9\), pero no tienen el mismo estatuto: \(\Pi\) publica
la única matriz APP y \(\Sigma\) genera sus acumulaciones. Sus versiones residuales
son
\[
 \bar\Sigma([a]_9,[b]_9)=[a+b]_9,
 \qquad
 \bar\Pi([a]_9,[b]_9)=[ab]_9.
\]
El cambio de carta es compatible con ambas operaciones:
\begin{align*}
 \lambda_9\circ\Sigma
   &=\bar\Sigma\circ\lambda_9^{\times2},\\
 \lambda_9\circ\Pi
   &=\bar\Pi\circ\lambda_9^{\times2}.
\end{align*}

Para no perder el cociente entero se introducen los levantamientos
\begin{align}
 a+b &= R^+_{a,b}+9Q^+_{a,b},
 &1\leq R^+_{a,b}\leq9,
\label{p12-83-c1-s2:eq:app-lift-suma}\\
 ab &= R^\times_{a,b}+9Q^\times_{a,b},
 &1\leq R^\times_{a,b}\leq9.
\label{p12-83-c1-s2:eq:app-lift-producto}
\end{align}
Aquí
\[
 R^+_{a,b}=\Sigma(a,b),
 \qquad
 R^\times_{a,b}=\Pi(a,b),
\]
y los cocientes están determinados unívocamente por
\[
 Q^+_{a,b}=\frac{a+b-R^+_{a,b}}9,
 \qquad
 Q^\times_{a,b}=\frac{ab-R^\times_{a,b}}9.
\]

\begin{definition}[APP genealógica]
\label{p12-83-c1-s2:def:app-genealogica-autonoma}
La estructura APP en esta etapa es la tupla
\[
 \operatorname{APP}:=
 \bigl(
  X_9,
  \Gamma_{\rm APP},
  \mathsf{Path}(\Gamma_{\rm APP}),
  \Sigma,\Pi,
  R^+,Q^+,R^\times,Q^\times
 \bigr).
\]
La genealogía de una trayectoria conserva, para cada vértice visitado, la
posición, la dirección de entrada, la entrada multiplicativa, la acumulación
aditiva y sus respectivos residuos y cocientes.
\end{definition}

\begin{theorem}[Construcción causal del soporte APP]
\label{p12-83-c1-s2:thm:app-cadena-constructiva-autonoma}
La cadena
\[
 \mathcal D_9
 \longrightarrow
 \mathcal D_8^2
 \longrightarrow
 \mathcal D_8^2\sqcup\mathcal G_9
 \longrightarrow
 \mathcal D_9^2
 \xrightarrow{\lambda_9^{\times2}}
 X_9
 \longrightarrow
 \Gamma_{\rm APP}
 \longrightarrow
 \mathsf{Path}(\Gamma_{\rm APP})
 \longrightarrow
 (\Sigma,\Pi)
\]
construye, sin pérdida cardinal, el soporte, el movimiento, las trayectorias,
la matriz multiplicativa de APP y su ley aditiva interna.
\end{theorem}

\begin{proof}
La \cref{p12-83-c1-s2:prop:app-minimalidad-completacion} construye la carta completa a
partir del núcleo. El \cref{p12-83-c1-s2:thm:app-equivalencia-modelos} identifica su
modelo intersticial. El \cref{p12-83-c1-s2:lem:app-cambio-carta-sin-perdida} transporta
los ochenta y un estados a \(X_9\). El grafo de Cayley incorpora las cuatro
traslaciones reversibles sin cambiar los vértices. La categoría libre
incorpora todas las concatenaciones finitas. Finalmente, \(\Sigma\) y
\(\Pi\) son aplicaciones totales distintas sobre el mismo dominio; las
ecuaciones \cref{p12-83-c1-s2:eq:app-lift-suma,p12-83-c1-s2:eq:app-lift-producto} conservan la parte
que la reducción residual no publica.
\end{proof}

\subsection{Descomposición del ortograma APP en núcleo \(8\times8\), gnomon de \(17\) estados y carta \(9\times9\) con lados opuestos identificados}

La \cref{fig:app-ortograma} debe leerse en tres capas: el cuadrado interior
de sesenta y cuatro estados, el gnomon de diecisiete estados y la
identificación de lados opuestos. La carta de nueve por nueve completa así el
núcleo de ocho por ocho mediante el gnomon de cierre. La matriz
multiplicativa y la ley aditiva interna se evalúan en cada vértice del mismo
estado, conservan levantamientos residuo--cociente distintos y no constituyen
dos cuadrados APP.

\subsection{Obstrucciones de fidelidad del soporte APP: cardinales \(81/17\), traslaciones reversibles y separación de las hojas suma y producto}

\begin{criterion}[Criterios de refutación del soporte APP]
La construcción queda refutada si ocurre cualquiera de los hechos
siguientes:
\begin{enumerate}
 \item el soporte completo tiene cardinal distinto de \(81\);
 \item el complemento gnomónico no tiene cardinal \(17\);
 \item un estado de \(X_9\) carece de alguna de las cuatro traslaciones;
 \item se identifican \(\Sigma\) y \(\Pi\) porque coincidan en un valor;
 \item se elimina \(Q^+\) o \(Q^\times\) al retornar por una frontera;
 \item se presentan las trayectorias de distinta longitud como nuevos
       vértices del soporte;
 \item se interpreta \(9\) y \([0]_9\) como fases independientes.
\end{enumerate}
\end{criterion}

La unidad ha construido APP hasta su evaluación genealógica completa. Todavía
no ha orientado la diferencia entre la matriz y su acumulador, no ha
introducido el TRIT
y no ha definido el Topological Phase Kernel. La unidad siguiente estudiará
la estructura algebraica de filas y columnas, el sector de unidades y el
cociclo de acarreo que prepara esa orientación.


% Unidad U005. Redacción autónoma desde 03c/03h/20 de la autoridad 1063/live.

\long\def\HMTDeferredTPKBase{%
\section{Topological Phase Kernel: base observable y estado enriquecido}
\label{p12-83-c1-s3:chap:tpk-categoria-estado-enriquecido}

APP ha proporcionado un soporte, una matriz multiplicativa y una ley aditiva
interna; TRIT ha
proporcionado una orientación levantada. El Topological Phase Kernel, en
adelante TPK, especifica cómo se transportan conjuntamente esos datos, qué
coordenadas se conservan al cerrar una ruta y qué lectores pueden publicarse
sin identificar la publicación con el estado. Su objeto matemático es una
categoría sobre una base de caminos, provista de una familia explícita de
fibras y relaciones de prolongación.

\subsection{La base doble de caminos APP}

Sea \(\mathsf P_{\rm APP}\) la categoría libre construida en la
\cref{p12-83-c1-s2:chap:app-completacion-gnomonica}. Definimos
\[
 \mathsf P_{\rm APP}^{(2)}
 :=\mathsf P_{\rm APP}\times\mathsf P_{\rm APP}.
\]
Un objeto es un par de posiciones
\((p^\Sigma,p^\Pi)\in X_9^2\). Un morfismo es un par de trayectorias; una
de ellas puede ser la identidad. La primera componente se evalúa por
\(\Sigma\), la segunda por \(\Pi\). No son dos retículas: son dos cursores
tipados sobre el mismo soporte APP.

El conjunto de direcciones cardinales es
\[
 \mathscr D_{\rm card}:=\{N,E,S,O\},
\]
con desplazamientos
\[
 \delta(N)=(-1,0),\quad
 \delta(E)=(0,1),\quad
 \delta(S)=(1,0),\quad
 \delta(O)=(0,-1)
 \quad\text{en }X_9.
\]
La involución de orientación es
\[
 \overline N=S,
 \quad \overline S=N,
 \quad \overline E=O,
 \quad \overline O=E.
\]

\subsection{Fase nonádica, sector dodecafásico y calendario}

El reloj completo de esta etapa es
\[
 \Theta_{108}:=\mathbb Z/108\mathbb Z.
\]
Para \(\theta\in\Theta_{108}\), sea
\(\widetilde\theta\in\{0,\ldots,107\}\) su representante normalizado.
Definimos tres coordenadas distintas:
\begin{align*}
 g_0(\theta)&=[\widetilde\theta+1]_9\in C_9,\\
 r(\theta)&=1+(\widetilde\theta\bmod9)\in\mathcal D_9,\\
 \beta(\theta)&=
 \left[\left\lfloor\frac{\widetilde\theta}{9}\right\rfloor\right]_{12}
 \in C_{12}.
\end{align*}
La primera es la clase de fase, la segunda su etiqueta positiva y la tercera
el sector dodecafásico. La aplicación \(\beta\) es una aplicación de bloques;
no es la reducción canónica \(\theta\bmod12\).

La firma y la hoja activa son
\[
 \tau(\theta):=\varepsilon_9(r(\theta))\in\mathsf{TRIT},
\]
y
\[
 \mu(\theta)=
 \begin{cases}
  \Sigma,&\tau(\theta)=+1,\\
  \Pi,&\tau(\theta)=-1,\\
  \varnothing,&\tau(\theta)=0.
 \end{cases}
\]
El símbolo \(\varnothing\) designa un paso de umbral en el que ningún cursor
se traslada; no designa ausencia de estado.

\begin{definition}[Configuración observable]
\label{p12-83-c1-s3:def:tpk-configuracion-observable}
La base observable es
\[
 \mathcal Q_{\rm obs}
 :=X_9^2\times\mathscr D_{\rm card}^2\times\Theta_{108}.
\]
Un elemento se escribe
\[
 q=(p^\Sigma,p^\Pi,d^\Sigma,d^\Pi,\theta).
\]
Contiene dos posiciones, dos orientaciones cardinales y un índice de
calendario.
\end{definition}

\subsection{Selector de hoja \(M_{\rm ph}:\mathcal D_9\to\mathsf{TRIT}\) y evolución conjunta de posición, orientación y calendario}

La función
\[
 M_{\rm ph}:\mathcal D_9\longrightarrow\mathsf{TRIT},
 \qquad M_{\rm ph}(r)=\varepsilon_9(r)
\]
es el selector interno de hoja. Selecciona qué evaluación actúa; no desplaza
por sí mismo ninguna posición.

Definimos el indicador de inversión de orientación por
\[
 \eta_{27}(\theta)=
 \begin{cases}
  -1,&\widetilde\theta+1\in\{27,54,81,108\},\\
  +1,&\text{en otro caso}.
 \end{cases}
\]
Las nuevas direcciones son
\[
 d^{\nu,+}=
 \begin{cases}
  \overline{d^\nu},&\eta_{27}(\theta)=-1,\\
  d^\nu,&\eta_{27}(\theta)=+1,
 \end{cases}
 \qquad \nu\in\{\Sigma,\Pi\}.
\]
Las posiciones evolucionan mediante
\[
 (p^{\Sigma,+},p^{\Pi,+})=
 \begin{cases}
 (p^\Sigma+\delta(d^\Sigma),p^\Pi),&\tau(\theta)=+1,\\
 (p^\Sigma,p^\Pi+\delta(d^\Pi)),&\tau(\theta)=-1,\\
 (p^\Sigma,p^\Pi),&\tau(\theta)=0.
 \end{cases}
\]

\begin{definition}[Evolución observable]
\label{p12-83-c1-s3:def:tpk-evolucion-observable}
La aplicación
\[
 F_{\rm obs}:\mathcal Q_{\rm obs}\longrightarrow\mathcal Q_{\rm obs}
\]
se define por
\[
 F_{\rm obs}(q)=
 (p^{\Sigma,+},p^{\Pi,+},d^{\Sigma,+},d^{\Pi,+},[\theta+1]_{108}).
\]
La categoría libre del grafo \(q\to F_{\rm obs}(q)\) se denota por
\(\mathsf P_{\rm obs}\).
\end{definition}

Cada flecha de \(\mathsf P_{\rm obs}\) induce un par de trayectorias en
APP, de modo que existe un funtor
\[
 B:\mathsf P_{\rm obs}\longrightarrow
 \mathsf P_{\rm APP}^{(2)}.
\]
Si \(\tau=+1\), la segunda componente de \(B\) es identidad; si
\(\tau=-1\), lo es la primera; si \(\tau=0\), ambas lo son.

\begin{proposition}[Separación entre selección, movimiento y observación]
\label{p12-83-c1-s3:prop:tpk-separacion-seleccion-movimiento}
El selector \(M_{\rm ph}\), las traslaciones
\(T_d(p)=p+\delta(d)\) y cualquier muestreo temporal poseen dominios y
codominios distintos. No son operadores intercambiables.
\end{proposition}

\begin{proof}
\(M_{\rm ph}\) toma una etiqueta de fase y devuelve una orientación trítica.
\(T_d\) es un automorfismo de \(X_9\). Un muestreo actúa sobre una
sucesión o una categoría de flechas. Identificar dos de ellos violaría el
tipado. \(F_{\rm obs}\) los compone en un orden declarado sin borrar esa
distinción.
\end{proof}

\subsection{Descomposición del estado enriquecido en \(6\) familias fibradas}

Sobre cada \(q\in\mathcal Q_{\rm obs}\) fijamos la fibra total
\[
 \mathcal F_q=
 \mathsf O_q\times
 \mathsf H_q\times
 \mathsf C_q\times
 \mathsf S_q\times
 \mathsf B_q\times
 \Lambda_q.
\]
Cada factor tiene un tipo y una función propios.

\subsubsection{Determinación conjunta de la orientación y la hoja del estado enriquecido}

\[
 \mathsf O_q\subseteq
 \mathscr D_{\rm card}^2\times\{+1,-1\}.
\]
Las dos direcciones coinciden con las almacenadas en \(q\); la última
coordenada registra la orientación de hoja cuando la realización la requiere.
La etiqueta de hoja no sustituye las direcciones cardinales.

\subsubsection{Descomposición euclídea del estado en residuo y cociente compatibles}

La primera fibra de elevación es
\[
 \mathsf H^{(1)}=\mathcal D_9\times\mathbb Z,
 \qquad
 \operatorname{rec}_9(r,u)=r+9u.
\]
Para cada entero \(n\) existen únicos
\[
 r=\operatorname{dr}_9(n),
 \qquad
 u=\frac{n-r}{9}.
\]
Una profundidad mayor consiste en una torre declarada de estas fibras y sus
mapas de truncamiento. El término ``Hensel'' no sustituye la definición de
esa torre.

\subsubsection{Operador de acarreo y actualización conjunta de residuo, cociente y memoria}

\(\mathsf C_q\) es un grupo abeliano. En los lectores de bloques decimales
se toma \(\mathsf C_q=\mathbb Z\). El incremento no es una coordenada
independiente, sino una cocadena sobre las flechas:
\[
 \kappa(r)\in\mathsf C_{q'}
 \quad\text{para }r:q\to q'.
\]

\subsubsection{Construcción de la firma entera a partir de orientación, acarreo y frontera}

La firma se deriva mediante una aplicación tipada
\[
 \operatorname{Sig}_q:
 \mathsf H_q\times\mathsf C_q\times\Lambda_q
 \longrightarrow\mathsf S_q.
\]
En la realización de frontera nonádica se conservan, de manera individual,
\[
 q_\partial,a_\partial,c_\partial\in\mathbb F_3^6,
 \qquad
 \operatorname{colw}\in\mathbb Z_{\geq0}^{6},
 \qquad
 \rho_W,r_\partial\in\mathbb F_3^3.
\]
Estas coordenadas son margen, combinación axial, condición cúbica, pesos de
columna, clase residual de Witt y retorno ternario. No se condensan en una
única palabra si una prueba posterior utiliza alguna de ellas.

\subsubsection{Cilindro y frontera}

\[
 \mathsf B_q:=\mathsf{Cyl}_q\times\mathsf{Fr}_q.
\]
El cilindro localiza una familia de prefijos; la frontera determina qué
extensiones son compatibles. La supervivencia será una relación entre
estados provistos de ambos datos, no un número añadido a posteriori.

\subsubsection{Registro genealógico de transporte}

\(\Lambda_q\) contiene los registros de transporte cuyo extremo es
\(q\). Una prolongación \(r:q\to q'\) actúa por concatenación:
\[
 \lambda\longmapsto r\circ\lambda.
\]
El registro genealógico conserva la ruta; su extremo observable no la determina.

\subsection{Transporte de fibras}

Para cada flecha \(r:q\to q'\) se declaran aplicaciones
\[
 \mathsf H(r):\mathsf H_q\to\mathsf H_{q'},
 \quad
 \mathsf C(r):\mathsf C_q\to\mathsf C_{q'},
 \quad
 \Lambda(r):\Lambda_q\to\Lambda_{q'}.
\]
Respetan identidades y composición, y \(\mathsf C(r)\) es un homomorfismo.
El acarreo satisface
\begin{equation}
 \kappa(r_2\circ r_1)
 =\kappa(r_2)+\mathsf C(r_2)\bigl(\kappa(r_1)\bigr).
\label{p12-83-c1-s3:eq:tpk-cocadena-transporte}
\end{equation}
Por tanto, el transporte de las coordenadas determinadas es
\begin{align*}
 h'&=\mathsf H(r)(h),\\
 c'&=\mathsf C(r)(c)+\kappa(r),\\
 \lambda'&=\Lambda(r)(\lambda),\\
 \sigma'&=\operatorname{Sig}_{q'}(h',c',\lambda').
\end{align*}

\begin{lemma}[Asociatividad del acarreo transportado]
\label{p12-83-c1-s3:lem:tpk-asociatividad-acarreo}
La ecuación \cref{p12-83-c1-s3:eq:tpk-cocadena-transporte} garantiza que transportar un
estado por \(r_1\) y después por \(r_2\) produce el mismo acarreo que
transportarlo por \(r_2\circ r_1\).
\end{lemma}

\begin{proof}
El transporte sucesivo da
\[
 \mathsf C(r_2)\bigl(\mathsf C(r_1)c+\kappa(r_1)\bigr)+\kappa(r_2).
\]
Por functorialidad de \(\mathsf C\), el primer término es
\(\mathsf C(r_2r_1)c\); la suma de los otros dos es
\(\kappa(r_2r_1)\) por \cref{p12-83-c1-s3:eq:tpk-cocadena-transporte}.
\end{proof}

\subsection{Relaciones de extensión}

La orientación y la frontera no se fuerzan mediante una función universal.
Para cada flecha \(r:q\to q'\), se declara una relación
\[
 \operatorname{Ext}_r\subseteq\mathcal F_q\times\mathcal F_{q'}.
\]
Se exigen las dos condiciones
\[
 \Delta_q\subseteq\operatorname{Ext}_{\operatorname{id}_q},
 \qquad
 \operatorname{Ext}_{r_2}\circ\operatorname{Ext}_{r_1}
 \subseteq\operatorname{Ext}_{r_2\circ r_1}.
\]
La primera conserva identidades; la segunda conserva composición. Una
relación puede admitir varias extensiones locales y dejar una sola tras un
horizonte mayor.

\begin{definition}[Topological Phase Kernel]
\label{p12-83-c1-s3:def:tpk-categoria-autonoma}
Un objeto del TPK es una tupla
\[
 x=(q,o,h,c,\sigma,C,b,\lambda)
\]
tal que
\[
 q\in\mathcal Q_{\rm obs},\quad
 o\in\mathsf O_q,\quad
 h\in\mathsf H_q,\quad
 c\in\mathsf C_q,\quad
 \lambda\in\Lambda_q,
\]
\[
 \sigma=\operatorname{Sig}_q(h,c,\lambda),
 \qquad
 (C,b)\in\mathsf{Cyl}_q\times\mathsf{Fr}_q.
\]
Asociamos a cada objeto su componente fibral
\[
 \mathbf f(x):=(o,h,c,\sigma,(C,b),\lambda)\in\mathcal F_q.
\]
Un morfismo \(x\to x'\) es una terna \((r,x,x')\) en la que
\(r:q\to q'\) pertenece a \(\mathsf P_{\rm obs}\), las coordenadas
transportables satisfacen las leyes anteriores y
\((\mathbf f(x),\mathbf f(x'))\in
\operatorname{Ext}_r\). La composición es
\[
 (r_2,x',x'')\circ(r_1,x,x')=(r_2r_1,x,x''),
\]
y la identidad de \(x\) es \((\operatorname{id}_q,x,x)\).
\end{definition}

\begin{theorem}[Estructura categorial del TPK]
\label{p12-83-c1-s3:thm:tpk-es-categoria}
La \cref{p12-83-c1-s3:def:tpk-categoria-autonoma} define una categoría y la aplicación
\[
 p:\mathsf{TPK}\longrightarrow\mathsf P_{\rm obs},
 \qquad p(x)=q,\qquad p(r,x,x')=r,
\]
es un funtor.
\end{theorem}

\begin{proof}
La composición de rutas es asociativa. El
\cref{p12-83-c1-s3:lem:tpk-asociatividad-acarreo} prueba la compatibilidad de la
coordenada de acarreo; la functorialidad declarada de \(\mathsf H\),
\(\mathsf C\) y \(\Lambda\) prueba la de las demás coordenadas
determinadas. La estabilidad de \(\operatorname{Ext}\) garantiza que el
compuesto sigue siendo admisible, y su condición diagonal proporciona las
identidades. La aplicación \(p\) conserva identidades y composición por
construcción.
\end{proof}

\subsection{Proyecciones y pérdidas tipadas}

Las proyecciones elementales son
\[
 \phi(x)=g_0(\theta),
 \qquad
 \tau(x)=\varepsilon_9(r(\theta)),
 \qquad
 \mu(x)=\mu(\theta).
\]
La proyección al estado TRIT local es
\[
 \pi_{\rm TRIT}(x)
 =\bigl(\tau,\mu,o,h,c,\sigma,\lambda,\phi\bigr).
\]
Olvida el cilindro, la frontera y parte de \(q\); no define otro contenedor.
Dos estados pueden tener la misma imagen por \(\pi_{\rm TRIT}\) y diferir
en la información olvidada.

\subsection{Criterio de supervivencia de rutas mediante conservación de la frontera proyectiva}

Sea \(\mathcal X_{\rm TPK}^{(n)}\) la familia de estados a profundidad
\(n\). Para \(N\geq n\), definimos
\[
 \operatorname{Surv}_n^{(N)}=
 \left\{x_n:
 \begin{array}{l}
 \text{existen }x_{n+1},\ldots,x_N\text{ con}\\
 (x_j,x_{j+1})\in\operatorname{Ext}_{r_j}
 \text{ para }n\leq j<N
 \end{array}\right\}.
\]

\begin{proposition}[Filtración de supervivencia]
\label{p12-83-c1-s3:prop:tpk-filtracion-supervivencia}
Para \(N\geq n\),
\[
 \operatorname{Surv}_n^{(N+1)}
 \subseteq
 \operatorname{Surv}_n^{(N)}.
\]
Si el esquema está definido a toda profundidad, los estados
indefinidamente prolongables son
\[
 \operatorname{Surv}_n
 =\bigcap_{N\geq n}\operatorname{Surv}_n^{(N)}.
\]
\end{proposition}

\begin{proof}
Una prolongación hasta \(N+1\) se trunca a una prolongación hasta \(N\),
lo que prueba la inclusión. La intersección expresa exactamente la
prolongabilidad para todo horizonte finito.
\end{proof}

\subsection{Árbol operatorio del contenedor}

\begin{figure}[htbp]
 \centering
 \begin{tikzpicture}[
  node distance=7mm and 9mm,
  every node/.style={font=\small,align=center},
  root/.style={draw=black,very thick,rounded corners,inner sep=5pt},
  op/.style={draw=black,rounded corners,inner sep=4pt},
  edge/.style={-{Latex[length=2mm]},semithick}
 ]
  \node[root] (tpk) {TPK\\$\mathcal X_{\rm TPK}^{\rm enr}$};
  \node[op,below left=of tpk] (sel) {selección\\$M_{\rm ph}$};
  \node[op,below=of tpk] (mov) {transporte\\$T_d,F_{\rm obs}$};
  \node[op,below right=of tpk] (fib) {fibras\\$\mathsf H,\mathsf C,\mathsf S,
    \mathsf B,\Lambda$};
  \node[op,below=17mm of mov] (cal) {calendario\\$C_{12}\hookrightarrow
    C_{108}\twoheadrightarrow C_9$};
  \node[op,below=of cal] (ext) {extensión, cilindros y supervivencia};
  \draw[edge] (tpk) -- (sel);
  \draw[edge] (tpk) -- (mov);
  \draw[edge] (tpk) -- (fib);
  \draw[edge] (sel) |- (cal);
  \draw[edge] (mov) -- (cal);
  \draw[edge] (fib) |- (cal);
  \draw[edge] (cal) -- (ext);
 \end{tikzpicture}
 \caption{Contención operatoria del Topological Phase Kernel. El selector,
 el transporte, los lectores y las fibras no son contenedores paralelos: son
 mapas o coordenadas del mismo estado enriquecido.}
 \label{p12-83-c1-s3:fig:tpk-arbol-operatorio-autonomo}
\end{figure}

\subsection{Obstrucciones del estado enriquecido del TPK: tipado del selector, cocadena de acarreo, extensión compatible y conservación genealógica}

\begin{criterion}[Criterios de refutación del estado enriquecido]
La definición queda refutada si:
\begin{enumerate}
 \item el selector \(M_{\rm ph}\) se usa como traslación;
 \item se identifica \(\theta\), \(g_0(\theta)\), \(r(\theta)\) y
       \(\beta(\theta)\);
 \item se permite que una firma se elija independientemente de
       \((h,c,\lambda)\);
 \item la composición de acarreos incumple
       \cref{p12-83-c1-s3:eq:tpk-cocadena-transporte};
 \item se sustituye \(\operatorname{Ext}_r\) por una elección funcional no
       demostrada;
 \item se llama estado completo a una proyección que ha olvidado cilindro,
       frontera o registro genealógico;
 \item el retorno de fase borra alguna de las coordenadas de fibra.
\end{enumerate}
\end{criterion}

Esta unidad ha definido el contenedor. La siguiente enumerará sus lectores,
operadores, odómetro y calendarios, y demostrará por qué los retornos de
nueve, veintisiete y ciento ocho pasos no son reinicios del estado.
}%


% Unidad U001. Redacción autónoma; no procede del borrador condensado.

\section{Segmentos marcados, intersticios y célula nonádica}
\label{p12-83-c1-s4:chap:segmentos-marcas-intersticios}

La construcción comienza antes de toda constante y antes de toda elección de
unidad física. El dato inaugural no es un número real ya constituido, sino un
segmento ordenado, una familia de marcas y la relación de adyacencia entre
marcas consecutivas. Esta separación es necesaria porque una marca, un
intervalo, la longitud de ese intervalo y el numeral que publica dicha
longitud son objetos de tipos diferentes.

\subsection{Carta posicional finita y proyección residual decimal}

\begin{definition}[Segmento marcado de orden diez]
Sea
\[
 \mathsf M_{10}:=\{m_0,m_1,\ldots,m_{10}\}
\]
un conjunto totalmente ordenado por
\(m_0<m_1<\cdots<m_{10}\). Su grafo de adyacencia es el camino
\[
 \mathsf P_{10}:\quad
 m_0\longrightarrow m_1\longrightarrow\cdots\longrightarrow m_{10}.
\]
Denotaremos por
\[
 \mathsf I_k:=(m_k,m_{k+1}),\qquad 0\leq k\leq9,
\]
el intersticio abstracto comprendido entre dos marcas consecutivas. La
familia de intersticios es
\(\mathsf E_{10}:=\{\mathsf I_0,\ldots,\mathsf I_9\}\).
\end{definition}

La notación abierta \((m_k,m_{k+1})\) no presupone aún la recta real. Sólo
indica el par ordenado de extremos y la arista que los une. Cuando se elige la
carta arquimediana
\[
 \chi_{10}(m_k)=k,
\]
el intersticio abstracto \(\mathsf I_k\) se publica mediante el intervalo
semiabierto \([k,k+1)\). Por tanto, el índice \(k\), la arista
\(\mathsf I_k\) y el intervalo \([k,k+1)\) no se identificarán en el texto:
el primero es una etiqueta, el segundo es un elemento combinatorio y el
tercero es su realización en una carta.

\begin{lemma}[Censo de marcas e intersticios]
\label{p12-83-c1-s4:lem:censo-marcas-intersticios}
El segmento \(\mathsf P_{10}\) contiene once marcas y diez intersticios. Más
generalmente, un camino con \(n+1\) marcas posee exactamente \(n\) aristas.
\end{lemma}

\begin{proof}
La aplicación
\[
 \{0,\ldots,n-1\}\longrightarrow E(\mathsf P_n),
 \qquad k\longmapsto(m_k,m_{k+1}),
\]
es biyectiva. La afirmación particular se obtiene con \(n=10\).
\end{proof}

El extremo \(m_0\) fija el origen de la carta y \(m_{10}\) fija su umbral.
Ninguno de los dos constituye por sí mismo un intersticio adicional. Esta
observación elemental impide tres confusiones posteriores: contar extremos
como celdas, confundir el cierre con un nuevo estado de fase y tratar el cero
como una décima fase independiente del ciclo nonádico.

\subsection{Extracción de la célula de \(9\) fases}

La carta decimal distingue el anclaje \([0,1)\) de la célula periódica. Las
nueve marcas positivas
\[
 \mathsf D_9:=\{1,2,3,4,5,6,7,8,9\}
\]
se utilizarán como representantes públicos de
\(\mathbb Z/9\mathbb Z\), haciendo corresponder la marca \(9\) con la clase
residual \(\overline 0\). Definimos
\[
 \lambda_9:\mathsf D_9\longrightarrow\mathbb Z/9\mathbb Z,
 \qquad \lambda_9(a)=\overline a.
\]

\begin{proposition}[Carta positiva de las clases nonádicas]
\label{p12-83-c1-s4:prop:carta-positiva-nonadica}
La aplicación \(\lambda_9\) es biyectiva. Su inversa es la sección positiva
\[
 s_+(\overline a)=
 \begin{cases}
 a,&1\leq a\leq8,\\
 9,&\overline a=\overline0.
 \end{cases}
\]
En particular, la escritura pública \(9\) y la escritura residual \(0\)
representan la misma clase, pero pertenecen a cartas distintas.
\end{proposition}

\begin{proof}
Las nueve clases de \(\mathbb Z/9\mathbb Z\) poseen representantes únicos en
\(\{0,1,\ldots,8\}\). Sustituir el representante \(0\) por \(9\) produce
otra transversal completa, por lo que \(\lambda_9\) es biyectiva y \(s_+\)
es su inversa.
\end{proof}

La célula de fase no se obtiene borrando arbitrariamente el origen. El tramo
\([0,1)\) fija la carta decimal desde la que se observa el ciclo; las nueve
fases son las nueve clases que se publican como \(1,\ldots,9\). En la carta
levantada, el paso que cierra la célula se representa por \([9,10)\); en la
carta residual, ese mismo paso retorna de \(\overline0\) a \(\overline1\).
El cierre y el cambio de representante no son la misma operación.

\subsection{Posición, altura y división euclídea}

\begin{definition}[Coordenadas nonádicas]
Para cada entero \(n\geq1\), definimos
\[
 \rho_9(n):=1+((n-1)\bmod9),
 \qquad
 q_9(n):=\left\lfloor\frac{n-1}{9}\right\rfloor.
\]
La primera coordenada se denomina posición de fase; la segunda, altura de
vuelta o cociente de memoria.
\end{definition}

\begin{theorem}[Descomposición nonádica única]
\label{p12-83-c1-s4:thm:descomposicion-nonadica-unica}
Todo entero \(n\geq1\) admite una única escritura
\[
 n=9q+\rho,
 \qquad q\in\mathbb Z_{\geq0},\quad \rho\in\mathsf D_9.
\]
En esa escritura \(q=q_9(n)\) y \(\rho=\rho_9(n)\).
\end{theorem}

\begin{proof}
Aplicamos la división euclídea a \(n-1\): existen únicos
\(q\geq0\) y \(r\in\{0,\ldots,8\}\) tales que
\(n-1=9q+r\). Tomando \(\rho=r+1\) se obtiene
\(n=9q+\rho\), con \(1\leq\rho\leq9\). La unicidad de \(q,r\) implica la
de \(q,\rho\).
\end{proof}

\begin{corollary}[Retorno de fase con incremento de memoria]
\label{p12-83-c1-s4:cor:retorno-fase-incremento-memoria}
Para todo \(n\geq1\),
\[
 \rho_9(n+9)=\rho_9(n),
 \qquad
 q_9(n+9)=q_9(n)+1.
\]
Por consiguiente, una vuelta completa conserva la fase visible y cambia el
estado levantado.
\end{corollary}

\begin{proof}
Si \(n=9q+\rho\), entonces \(n+9=9(q+1)+\rho\); la unicidad del
\cref{p12-83-c1-s4:thm:descomposicion-nonadica-unica} da las dos identidades.
\end{proof}

La evolución por un paso queda escrita sin ambigüedad como
\[
 (q,\rho)\longmapsto
 \begin{cases}
 (q,\rho+1),&1\leq\rho<9,\\
 (q+1,1),&\rho=9.
 \end{cases}
\]
La segunda línea no es un reinicio. Recupera la etiqueta inicial de fase y
aumenta el cociente que registra cuántas vueltas se han realizado.

\subsection{Periodicidad de la proyección circular e inyectividad del levantamiento helicoidal nonádico}

Sea
\[
 \vartheta_9(\rho)=\frac{2\pi(\rho-1)}9.
\]
La proyección de fase sobre el círculo unitario es
\[
 \mathcal R_9(n)=
 \bigl(\cos\vartheta_9(\rho_9(n)),
       \sin\vartheta_9(\rho_9(n))\bigr).
\]
Para conservar simultáneamente la posición y la altura introducimos la
inmersión helicoidal
\[
 \mathcal H_9(n)=
 \left(
  \cos\vartheta_9(\rho_9(n)),
  \sin\vartheta_9(\rho_9(n)),
  q_9(n)+\frac{\rho_9(n)-1}{9}
 \right).
\]

\begin{proposition}[La proyección circular no determina el estado]
\label{p12-83-c1-s4:prop:rueda-no-determina-estado}
La aplicación \(\mathcal R_9\) es periódica de periodo nueve. La aplicación
\(\mathcal H_9\) es inyectiva. En particular,
\[
 \mathcal R_9(n+9)=\mathcal R_9(n),
 \qquad
 \mathcal H_9(n+9)\neq\mathcal H_9(n).
\]
\end{proposition}

\begin{proof}
La primera igualdad se sigue del
\cref{p12-83-c1-s4:cor:retorno-fase-incremento-memoria}. La tercera coordenada de
\(\mathcal H_9(n+9)\) supera en una unidad la de \(\mathcal H_9(n)\), por
lo que ambos puntos son distintos. Finalmente, de la tercera coordenada se
recupera \(n\) multiplicando por nueve y añadiendo uno a la parte de fase;
por tanto, \(\mathcal H_9\) es inyectiva.
\end{proof}

Esta proposición es el modelo mínimo del principio que recorrerá toda la
obra:
\[
 \boxed{\text{retorno de fase}\neq\text{retorno del estado}.}
\]
La conexión nonádica del Topological Phase Kernel conservará más campos que
la hélice elemental ---hoja, orientación, acarreo, frontera, terminal y
memoria de ruta---, pero no alterará esta distinción inaugural.

\subsection{Correspondencia entre cardinalidad discreta y medida proyectiva}

\begin{definition}[Conteo]
El conteo de un conjunto finito \(X\) es su cardinal \(|X|\). No requiere
una orientación, una escala ni un lector adicional.
\end{definition}

\begin{definition}[Dato de medida]
Un dato de medida discreta es una terna
\[
 (x,\mathcal L,\mathcal C),
\]
donde \(x\) es un estado, \(\mathcal L\) es un lector tipado y
\(\mathcal C\) es una carta de publicación compatible con el codominio de
\(\mathcal L\). Su salida visible es
\[
 \mathcal C(\mathcal L(x)).
\]
\end{definition}

El mismo conjunto puede tener un único cardinal y admitir múltiples medidas,
porque pueden variar la cocadena local, la orientación de integración o la
carta de publicación. Inversamente, dos estados distintos pueden publicar la
misma coordenada. Por ello la obra conservará siempre dos niveles:
\[
 \text{estado con genealogía}
 \quad\longrightarrow\quad
 \text{coordenada publicada}.
\]

\begin{criterion}[Criterios de refutación de la unidad inaugural]
La construcción queda refutada si ocurre cualquiera de los siguientes
hechos:
\begin{enumerate}
 \item se cuentan diez fases independientes en la célula residual;
 \item se identifica la marca \(9\) con el entero \(0\) sin declarar el
 cambio de carta;
 \item el paso \(9\to1\) fuerza \(q\mapsto0\);
 \item dos enteros distintos poseen la misma pareja \((q_9,\rho_9)\);
 \item se atribuye una unidad física a la célula antes de definir el
 transductor dimensional.
\end{enumerate}
\end{criterion}

\subsection{Dominio inaugural de la célula nonádica y exclusión causal de APP, TRIT, TPK y constantes analíticas}

La unidad construida hasta aquí contiene únicamente:
\[
 \mathsf P_{10},\quad
 \mathsf D_9,\quad
 \lambda_9,\quad
 (q_9,\rho_9),\quad
 \mathcal R_9,\quad
 \mathcal H_9.
\]
No contiene todavía APP, TRIT, el Topological Phase Kernel ni una constante
analítica. La siguiente unidad tomará dos copias de la célula nonádica,
construirá su completación gnomónica y demostrará cómo aparece el soporte
bidimensional sobre el que actúan la matriz multiplicativa de APP y su ley
aditiva interna.


% Unidad U003B. Esqueleto ciclotómico de APP y completación binomial.
% Redacción autónoma a partir de los objetos, no de la arquitectura del PDF abreviado.

\section{Esqueleto ciclotómico de APP y completación binomial de la unidad}
\label{p12-83-c1-s5:chap:app-esqueleto-completacion}

El defecto entre la matriz APP y su acumulador interno no es un número
aislado. Su ley en
longitud cartesiana, su descomposición en órbitas de orden tres y la
completación de la celda nonádica producen tres objetos diferentes que deben
mantenerse tipados:
\[
 \text{defecto escalar}\quad D_d,
 \qquad
 \text{módulo ciclotómico}\quad W_{\rm APP},
 \qquad
 \text{completación}\quad\overline{\mathcal C}_9.
\]
El primero cuenta; el segundo conserva orientación y rango; el tercero añade
estratos de frontera. Las igualdades entre sus cardinales no autorizan a
identificar sus elementos.

\subsection{Ley del defecto a módulo fijo y longitud variable}

Para cada entero \(d\geq1\), sea
\[
 B_d:=\{1,\ldots,9\}^d.
\]
La raíz digital nonádica se representa por
\[
 \operatorname{dr}_9(n)=
 \begin{cases}
 n\bmod9,&n\not\equiv0\pmod9,\\
 9,&n\equiv0\pmod9.
 \end{cases}
\]
Definimos la acumulación aditiva y la evaluación matricial multiplicativa
sobre el mismo dominio:
\[
 A_d(x_1,\ldots,x_d)
 =\operatorname{dr}_9\!\left(\sum_{j=1}^d x_j\right),
\]
\[
 M_d(x_1,\ldots,x_d)
 =\operatorname{dr}_9\!\left(\prod_{j=1}^d x_j\right).
\]
Sus masas visibles y su defecto son
\[
 S_\Sigma(d):=\sum_{x\in B_d}A_d(x),
 \qquad
 S_\Pi(d):=\sum_{x\in B_d}M_d(x),
\]
\[
 D_d:=S_\Pi(d)-S_\Sigma(d).
\]
Aquí \(d\) cuenta posiciones cartesianas. No varía el módulo \(9\) y no se
ha introducido todavía una dimensión física.

\begin{lemma}[Masa del acumulador aditivo]
\label{p12-83-c1-s5:lem:masa-hoja-aditiva-longitud}
Para todo \(d\geq1\),
\[
 S_\Sigma(d)=45\,9^{d-1}.
\]
\end{lemma}

\begin{proof}
Fijadas las primeras \(d-1\) coordenadas y un residuo publicado
\(r\in\{1,\ldots,9\}\), existe una única última coordenada que hace que la
suma tenga residuo \(r\). Cada valor visible posee así \(9^{d-1}\)
preimágenes. Sumando los pesos \(1+\cdots+9=45\) se obtiene la fórmula.
\end{proof}

\begin{lemma}[Censo de la matriz multiplicativa]
\label{p12-83-c1-s5:lem:censo-hoja-multiplicativa-longitud}
Para cada uno de los seis residuos unitarios hay \(6^{d-1}\) preimágenes
por valor de producto fijado. Los residuos \(3\) y \(6\) tienen
\(d6^{d-1}\) preimágenes cada uno. El residuo nulo, publicado como \(9\),
tiene
\[
 9^d-6^d-2d6^{d-1}
\]
preimágenes. En consecuencia,
\[
 S_\Pi(d)=9^{d+1}-9(d+3)6^{d-1}.
\]
\end{lemma}

\begin{proof}
El grupo de unidades de \(\mathbb Z/9\mathbb Z\) tiene seis elementos. Al
fijar \(d-1\) unidades, la última queda determinada por el producto
unitario prescrito, lo que da \(6^{d-1}\) soluciones para cada una de las
seis clases.

Para obtener producto de valoración ternaria exactamente uno, se elige cuál
de las \(d\) coordenadas pertenece a una de las dos clases \(3,6\), mientras
las demás son unidades. Fijado el producto \(3\) o \(6\), la última unidad
queda determinada; de ahí \(d6^{d-1}\) para cada clase. Todas las tuplas
restantes tienen producto nulo módulo \(9\). Restar de \(9^d\) los dos
censos anteriores da su multiplicidad. La suma ponderada por los
representantes visibles produce la última identidad.
\end{proof}

\begin{theorem}[Ley dimensional interna de APP]
\label{p12-83-c1-s5:thm:ley-dimensional-app-autonoma}
Para todo \(d\geq1\),
\[
 \boxed{D_d=4\,9^d-9(d+3)6^{d-1}.}
\]
En particular,
\[
 D_1=0,
 \qquad
 D_2=54,
\]
y
\[
 \sum_{d\geq1}D_dz^d
 =\frac{54z^2(1-3z)}{(1-9z)(1-6z)^2}.
\]
Además,
\[
 \frac{D_d}{9^d}
 =4-(d+3)\left(\frac23\right)^{d-1}
 \longrightarrow4.
\]
\end{theorem}

\begin{proof}
Se restan las fórmulas de los dos lemas. Para la función generatriz se usan
\[
 \sum_{d\geq1}9^dz^d=\frac{9z}{1-9z},
 \qquad
 \sum_{d\geq1}(d+3)6^{d-1}z^d
 =\frac{z(4-18z)}{(1-6z)^2},
\]
y se reduce el numerador común. La fórmula normalizada resulta de dividir
por \(9^d\); el término exponencial multiplicado por \(d+3\) converge a
cero.
\end{proof}

La igualdad \(D_2=54\) reproduce el censo \(459-405\), pero la fórmula
anterior y la ley de módulo variable
\(\Delta_{3^m}^{\rm pl}=m3^{2m-1}\) son enunciados distintos: la primera
varía \(d\) manteniendo el módulo \(9\); la segunda varía el módulo.

\subsection{Operador ciclotómico primario}

Sea
\[
 G_{A_2}=\begin{pmatrix}2&-1\\-1&2\end{pmatrix},
 \qquad
 C=\begin{pmatrix}0&-1\\1&-1\end{pmatrix}.
\]

\begin{proposition}[Coxeter de tipo \(A_2\)]
\label{p12-83-c1-s5:prop:coxeter-a2-autonomo}
Se verifican
\[
 C^3=I_2,
 \qquad
 C^{\mathsf T}G_{A_2}C=G_{A_2},
 \qquad
 \det(XI_2-C)=X^2+X+1.
\]
Por tanto \(C\) preserva la forma \(A_2\), tiene orden exactamente tres y
no posee vector fijo no nulo.
\end{proposition}

\begin{proof}
La multiplicación da
\[
 C^2=\begin{pmatrix}-1&1\\-1&0\end{pmatrix},
 \qquad C^3=I_2.
\]
La sustitución directa en \(C^{\mathsf T}G_{A_2}C\) devuelve
\(G_{A_2}\). Finalmente, \(1\) no es raíz de \(X^2+X+1\), luego
\(\ker(C-I_2)=0\).
\end{proof}

\subsection{Descomposición de APP en \(12\) fibras de tipo \(A_2\)}

Considérese la acción
\[
 a:\mathbb Z/9\mathbb Z\longrightarrow\mathbb Z/9\mathbb Z,
 \qquad a(r)=4r.
\]
Como \(4^3\equiv1\pmod9\), su restricción a las unidades tiene dos órbitas
libres:
\[
 \mathcal O_1=(1,4,7),
 \qquad
 \mathcal O_2=(2,8,5).
\]
Para una órbita \(\mathcal O=(r,4r,7r)\), definimos el módulo de aumento
\[
 A(\mathcal O):=
 \ker\!\left(\mathbb Z[\mathcal O]
 \xrightarrow{\,\sum\,}\mathbb Z\right).
\]
En la base
\[
 b_1=e_r-e_{4r},
 \qquad b_2=e_{4r}-e_{7r},
\]
la forma euclídea restringida tiene Gram \(G_{A_2}\), y el ciclo
\(r\mapsto4r\mapsto7r\mapsto r\) actúa por \(C\).

Para \(x=(x_1,x_2,x_3)\in(\mathbb Z/9\mathbb Z)^3\), los tres pares no
ordenados son
\[
 \{1,2\},\quad\{2,3\},\quad\{3,1\}.
\]
Sobre cada par se evalúan dos leyes tipadas:
\[
 \ell_{ij}^{\Sigma}(x)=x_i+x_j,
 \qquad
 \ell_{ij}^{\Pi}(x)=x_ix_j.
\]
Como
\[
 \ell_{ij}^{\Sigma}(ax)=a\ell_{ij}^{\Sigma}(x),
 \qquad
 \ell_{ij}^{\Pi}(ax)=a^{-1}\ell_{ij}^{\Pi}(x),
\]
la suma porta la orientación \(C\) y el producto la orientación inversa
\(C^{-1}\).

\begin{definition}[Módulo ciclotómico de APP]
\label{p12-83-c1-s5:def:modulo-ciclotomico-app-autonomo}
Se define
\[
 \boxed{
 W_{\rm APP}:=
 \bigoplus_{\{i,j\}}
 \bigoplus_{\mu\in\{\Sigma,\Pi\}}
 \bigoplus_{\mathcal O\in\{\mathcal O_1,\mathcal O_2\}}
 A(\mathcal O).}
\]
Hay \(3\cdot2\cdot2=12\) sumandos, cada uno de rango dos; por tanto
\[
 W_{\rm APP}\cong A_2^{12},
 \qquad \operatorname{rank}W_{\rm APP}=24.
\]
Las etiquetas par--ley--órbita son parte del objeto.
\end{definition}

Para \(u\in\mathbb Z/9\mathbb Z\), escribimos
\[
 \beta_{\mathcal O}(u)=
 \begin{cases}
 (1,0),&u=r,\\
 (0,1),&u=4r,\\
 (-1,-1),&u=7r,\\
 (0,0),&u\notin\mathcal O,
 \end{cases}
\]
en la base \((b_1,b_2)\). El mapa estatal es
\[
 \Psi(x)_{ij,\mu,\mathcal O}
 :=\beta_{\mathcal O}(\ell_{ij}^{\mu}(x)).
\]

\begin{theorem}[Equivarianza y rango del esqueleto estatal]
\label{p12-83-c1-s5:thm:esqueleto-estatal-app-autonomo}
Sea \(\rho(a)\) la acción que usa \(C\) en las seis fibras aditivas y
\(C^{-1}\) en las seis multiplicativas. Entonces
\[
 \boxed{\Psi(ax)=\rho(a)\Psi(x)}
\]
para los \(729\) estados de \((\mathbb Z/9\mathbb Z)^3\). El submódulo
racional generado por las imágenes tiene rango \(24\).
\end{theorem}

\begin{proof}
En la ley aditiva, multiplicar todas las coordenadas por \(4\) avanza una
posición en cada órbita; en la ley multiplicativa multiplica el valor por
\(4^2\equiv7\equiv4^{-1}\pmod9\), y por ello invierte el ciclo. Esto prueba
la equivarianza componente a componente.

Para que el rango no quede oculto en una afirmación informática, fijamos el
orden de coordenadas
\[
 (12,\Sigma,\mathcal O_1),(12,\Sigma,\mathcal O_2),
 (12,\Pi,\mathcal O_1),(12,\Pi,\mathcal O_2),
\]
seguido del mismo orden para los pares \(23\) y \(31\), y dentro de cada
fibra las dos coordenadas \((b_1,b_2)\). Consideremos los estados
\[
\begin{gathered}
001,002,004,005,010,011,012,013,014,015,016,020,\\
022,023,040,050,101,102,104,105,110,120,140,150.
\end{gathered}
\]
La matriz entera \(M_\Psi\) cuyas filas son sus veinticuatro imágenes tiene
entradas en ({-1,0,1}). La eliminación fraccionariamente libre de
Bareiss produce como menores principales sucesivos
\[
1,1,1,-1,-1,-1,-1,1,-1,1,1,1,1,-2,-1,1,1,1,2,1,4,-6,12,9.
\]
En particular,
\[
 \det M_\Psi=9\neq0.
\]
Estas veinticuatro imágenes son linealmente independientes. Como el
codominio ya tiene rango \(24\), generan \(W_{\rm APP}\otimes\mathbb Q\).
\end{proof}

La prueba exhibe un certificado finito reproducible; no deduce rango pleno
de la mera sobreyectividad de cada proyección componente.

\subsection{Localización del agregado \texorpdfstring{\(54\)}{54}}

Sea
\[
 \mathcal M:=\mathbb Z[\mathbb Z/9\mathbb Z]
\]
y, para \(\mu\in\{\Sigma,\Pi\}\),
\[
 \nu_\mu:=\sum_{x,y\in\mathbb Z/9\mathbb Z}
 e_{\ell^\mu(x,y)}.
\]
El censo de las dos tablas da
\[
 \delta_{\Pi\Sigma}:=\nu_\Pi-\nu_\Sigma
 =12e_0+3e_3+3e_6
 -3\sum_{u\in(\mathbb Z/9\mathbb Z)^\times}e_u.
\]
Sus coeficientes son constantes en las órbitas de \(a\); por tanto
\[
 (I-a)\delta_{\Pi\Sigma}=0.
\]
Definimos la proyección de aumento
\[
 T:\mathcal M\longrightarrow
 A(\mathcal O_1)\oplus A(\mathcal O_2),
\]
\[
 T(\nu)=
 \left(((I-a)\nu)|_{\mathcal O_1},
       ((I-a)\nu)|_{\mathcal O_2}\right).
\]
Entonces
\[
 T(\delta_{\Pi\Sigma})=0.
\]
Sin embargo, el lector visible
\[
 w(e_0)=9,
 \qquad w(e_r)=r\quad(1\leq r\leq8)
\]
satisface
\[
 w(\delta_{\Pi\Sigma})=54.
\]

\begin{corollary}[Obstrucción isotípica]
\label{p12-83-c1-s5:cor:obstruccion-isotipica-app-autonoma}
El escalar \(54\) pertenece al sector \(C_3\)-invariante y no puede
recuperarse mediante un funcional lineal que factorice exclusivamente por
los módulos de aumento. En particular,
\[
 \sum_x\Psi(x)=0
\]
y, para todo funcional lineal \(L\) sobre \(W_{\rm APP}\otimes\mathbb Q\),
\[
 L\!\left(\sum_x\Psi(x)\right)=0\neq54.
\]
\end{corollary}

Esta obstrucción impide presentar \(54\) como una suma lineal retrospectiva
de las doce fibras. El acoplamiento APP--Fock posterior conserva la
orientación y utiliza un objeto adicional; no contradice el corolario.

\subsection{Completación cúbica de la celda nonádica}

Sea \(E_9\) el conjunto subyacente de \(\mathbb Z/9\mathbb Z\), y sea
\(\star\notin E_9\) un estado de frontera. Definimos
\[
 \mathcal C_9:=E_9^3,
 \qquad
 \overline{\mathcal C}_9:=(E_9\sqcup\{\star\})^3.
\]
El símbolo \(\star\) no es la clase cero: indica que una coordenada ha
alcanzado el estrato de cierre.

\begin{theorem}[Estratificación binomial de la unidad cúbica]
\label{p12-83-c1-s5:thm:emrh-binomial-autonoma}
El estrato con exactamente \(r\) coordenadas iguales a \(\star\) contiene
\[
 \binom3r9^{3-r}
\]
elementos. Por tanto,
\[
 \boxed{1000=729+243+27+1=729+271.}
\]
Al retirar sólo el ápice \((\star,\star,\star)\), queda
\[
 999=729+243+27=729+270.
\]
\end{theorem}

\begin{proof}
Se eligen las \(r\) coordenadas de frontera de \(\binom3r\) maneras; cada
coordenada restante admite nueve valores. Los cuatro estratos son disjuntos
y agotan el producto cartesiano. La suma es la expansión de
\((9+1)^3\).
\end{proof}

\begin{figure}[htbp]
\centering
\resizebox{0.96\linewidth}{!}{\input{figuras/base_autoridad/fig_cubo_emrh_autoridad}}
\caption{Completación cúbica de la celda nonádica. El interior de
cardinal \(729\) y los estratos de frontera de cardinales \(243\), \(27\) y
\(1\) son partes disjuntas de un único objeto.}
\label{p12-83-c1-s5:fig:cubo-emrh-autonomo}
\end{figure}

Esta estratificación recibe en HMT el nombre de extrema y media razón
holográfica. Matemáticamente es la completación binomial tipada de la celda
nonádica; el nombre no sustituye su definición.

Para todo \(d\geq1\), sea
\[
 \overline{\mathcal C}^{(d)}_9=(E_9\sqcup\{\star\})^d.
\]
Asignemos peso \(1/10\) a cada símbolo de una coordenada y tomemos la medida
producto \(\mu_d\).

\begin{theorem}[Compatibilidad proyectiva de la medida]
\label{p12-83-c1-s5:thm:medida-proyectiva-completacion}
La proyección que olvida la última coordenada satisface
\[
 (\pi_d)_*\mu_d=\mu_{d-1}.
\]
Además,
\[
 10^d=\sum_{r=0}^{d}\binom dr9^{d-r},
 \qquad
 \mu_d(\overline{\mathcal C}^{(d)}_9)=1.
\]
El número de estados aumenta con \(d\), mientras la masa total normalizada
se conserva.
\end{theorem}

\begin{proof}
Para todo cilindro \(A\times(E_9\sqcup\{\star\})\),
\[
 \mu_d(A\times(E_9\sqcup\{\star\}))
 =\mu_{d-1}(A)\sum_{u\in E_9\sqcup\{\star\}}\frac1{10}
 =\mu_{d-1}(A).
\]
Los cilindros generan el álgebra finita de subconjuntos, lo que determina la
medida imagen. La identidad cardinal vuelve a ser el binomio de Newton.
\end{proof}

\subsection{Morfismo del bloque central de APP hacia su realización euclídea}

El bloque central ya construido es
\[
 B_c=\begin{pmatrix}7&2\\2&7\end{pmatrix}.
\]
La lectura por filas y la lectura diagonal dan
\[
 72+27=77+22=99,
\]
\[
 72-27=45=\det B_c,
 \qquad
 77-22=55=54+1.
\]
La relación con la completación no se obtiene concatenando símbolos, sino
por la división euclídea
\[
 729=10\cdot72+9,
 \qquad
 271=10\cdot27+1.
\]

\begin{theorem}[Unicidad del levantamiento decimal]
\label{p12-83-c1-s5:thm:levantamiento-decimal-emrh-autonomo}
Para \(L_d(n)=10n+d\), con \(d\in\{0,\ldots,9\}\), existe un único par
\((d,d')\) tal que
\[
 L_d(72)=729,
 \qquad
 L_d(72)+L_{d'}(27)=1000,
 \qquad
 d+d'=10.
\]
Es el par
\[
 (d,d')=(9,1).
\]
\end{theorem}

\begin{proof}
La primera igualdad es \(720+d=729\), luego \(d=9\). La segunda queda
\(729+270+d'=1000\), de donde \(d'=1\). La tercera se verifica y la
unicidad de cociente y resto excluye otra solución.
\end{proof}

Así se obtiene el diagrama causal
\[
 B_c\longrightarrow(72,27)
 \longleftarrow((72,9),(27,1))
 \longleftarrow(729,271).
\]
El bloque y su lectura pertenecen a APP; los dividendos y restos pertenecen
a la completación. El diagrama comparte cocientes, no invierte la
causalidad.

\subsection{Incidencia del cubo y separación respecto de hexadas y octadas}

Sea \(M_d\) la matriz de incidencia de las \(2d\) caras de codimensión uno
de un hipercubo sobre sus \(2^d\) vértices. Cada cara contiene
\(2^{d-1}\) vértices; dos caras opuestas son disjuntas y dos caras
transversales comparten \(2^{d-2}\) vértices.

\begin{theorem}[Espectro de incidencia hipercúbica]
\label{p12-83-c1-s5:thm:espectro-incidencia-cubo-autonomo}
\[
 \operatorname{Spec}(M_dM_d^{\mathsf T})
 =\left\{
 d2^{d-1}\,^{(1)},
 2^{d-1}\,^{(d)},
 0^{(d-1)}
 \right\}.
\]
Para \(d=3\),
\[
 \operatorname{Spec}(M_3M_3^{\mathsf T})
 =\{12,4,4,4,0,0\}.
\]
\end{theorem}

\begin{proof}
El espacio de filas se descompone en: la recta constante; las \(d\)
diferencias de cada par de caras opuestas; y las (d-1) combinaciones de
sumas de pares cuya suma total es cero. Las intersecciones descritas antes
hacen que el Gram actúe por \(d2^{d-1}\), \(2^{d-1}\) y \(0\),
respectivamente.
\end{proof}

El cubo tridimensional posee seis caras, ocho vértices y doce aristas. Esos
cardinales anticipan un patrón de incidencia (6/8), pero una cara no es
una hexada y el conjunto de vértices no es una octada. La identificación
excepcional exige un transportador posterior que conserve soportes y
calibre.

\subsection{Resonancia del operador de incidencia en rango \(6\)}

Definamos los estratos que conservan exactamente \(k\) coordenadas
nonádicas en rango \(d\):
\[
 \mathcal I_k(d):=\binom dk9^k.
\]
Los defectos de APP son
\[
 D_{\rm vis}=54,
 \qquad
 D_{\rm tot}=1215,
 \qquad
 D_{\rm coc}=129,
 \qquad
 1215=54+9\cdot129.
\]

\begin{theorem}[Unicidad de la resonancia binomial]
\label{p12-83-c1-s5:thm:resonancia-binomial-autonoma}
El único entero positivo \(d\) que satisface simultáneamente
\[
 \mathcal I_1(d)=54,
 \qquad
 \mathcal I_2(d)=1215
\]
es \(d=6\). En ese rango,
\[
 \frac{\mathcal I_2(6)-\mathcal I_1(6)}9=129.
\]
\end{theorem}

\begin{proof}
La primera igualdad es \(9d=54\), que fuerza \(d=6\). Entonces
\[
 \mathcal I_2(6)=81\binom62=81\cdot15=1215,
\]
y
\[
 \frac{1215-54}{9}=129.
\]
No queda otro rango posible porque la primera igualdad ya determina \(d\).
\end{proof}

\subsection{Obstrucciones del esqueleto ciclotómico de APP: fibras \(A_2^{12}\), completación \(729+271\), resonancia de rango \(6\) e incidencia cúbica}

\begin{criterion}[Falsación tipada de la unidad]
La construcción queda refutada si ocurre cualquiera de los hechos
siguientes:
\begin{enumerate}
 \item la ley de \(D_d\) no reproduce \(D_1=0\) y \(D_2=54\);
 \item se fusiona la variación de longitud con la variación de módulo;
 \item \(C\) deja de preservar \(G_{A_2}\) o adquiere un vector fijo;
 \item las etiquetas par--hoja--órbita no producen doce fibras;
 \item el menor certificado de \(\Psi\) tiene determinante nulo;
 \item un funcional que factoriza por el módulo de aumento recupera \(54\);
 \item el símbolo de frontera \(\star\) se identifica con el residuo cero;
 \item los cuatro estratos no suman \(1000\);
 \item se identifica una cara cúbica con una hexada sin un transportador de
       incidencia.
\end{enumerate}
\end{criterion}

La unidad ha construido, sin reconocerlos retrospectivamente, el módulo
\(A_2^{12}\), la completación \(729+271\), la medida proyectiva y la
resonancia de rango seis. Estas estructuras alimentarán de manera distinta
la generación de constantes y las ramas excepcionales.


% Unidad U006F. Inventario operatorio completo del Topological Phase Kernel.

\long\def\HMTDeferredTPKDevelopment{%
\section{Álgebra de operadores del Topological Phase Kernel y dinámica causal de actualización del estado enriquecido}
\label{p12-83-c1-s6:chap:inventario-operatorio-completo-tpk}

El Topological Phase Kernel no es una reducción módulo tres ni un reloj de
ciento ocho posiciones. Es una categoría de estados enriquecidos equipada
con operadores de selección, transporte, lectura, memoria, periodicidad,
prolongación y publicación. Para impedir que una abreviatura sustituya al
objeto, cada operador se registra por su tipo y por la información que
conserva.

\subsection{Esquema tipado de las 34 construcciones canónicas del TPK: dominio, codominio, acción, invariantes y memoria}

\begin{definition}[Ficha operatoria]
\label{p12-83-c1-s6:def:u006f-ficha-operatoria}
Una realización canónica del TPK es una séxtupla
\begin{equation}
 \mathcal O=
 (D_{\mathcal O},C_{\mathcal O},f_{\mathcal O},
 I_{\mathcal O},L_{\mathcal O},M_{\mathcal O}),
 \label{p12-83-c1-s6:eq:u006f-ficha}
\end{equation}
donde:
\begin{enumerate}
 \item \(f_{\mathcal O}:D_{\mathcal O}\to C_{\mathcal O}\) es su regla;
 \item \(I_{\mathcal O}\) es la familia de relaciones conservadas;
 \item \(L_{\mathcal O}\) declara la información publicada o perdida;
 \item \(M_{\mathcal O}\) enumera la memoria que permanece en el estado
 enriquecido.
\end{enumerate}
\end{definition}

Dos realizaciones representan el mismo operador sólo cuando existe una
equivalencia tipada que entrelaza reglas e invariantes. Compartir periodo,
cardinal o salida visible no basta.

\subsection{Clasificación del estado enriquecido en \(8\) clases estructurales invariantes}

El inventario se particiona en:
\begin{enumerate}
 \item[\textbf{C1.}] soportes y estados;
 \item[\textbf{C2.}] selección interna;
 \item[\textbf{C3.}] transporte;
 \item[\textbf{C4.}] lectores y cocientes;
 \item[\textbf{C5.}] memoria y frontera;
 \item[\textbf{C6.}] periodicidad y retorno;
 \item[\textbf{C7.}] prolongación;
 \item[\textbf{C8.}] salidas.
\end{enumerate}
La clase se determina por dominio, codominio y función, no por el nombre de
una cifra asociada.

\Needspace{.92\textheight}
\subsection{Descomposición funcional en \(14\) agrupaciones compatibles con los operadores APP–TRIT–TPK}

\begingroup
\small
\renewcommand{\arraystretch}{1.12}
\begin{longtable}{@{}>{\raggedright\arraybackslash}p{.20\textwidth}>{\raggedright\arraybackslash}p{.29\textwidth}>{\raggedright\arraybackslash}p{.43\textwidth}@{}}
\toprule
Agrupación&Operadores o espacios&Función e invariantes\\
\midrule
\endfirsthead
\toprule
Agrupación&Operadores o espacios&Función e invariantes\\
\midrule
\endhead
Estado y memoria de ruta&
\(\mathcal X_{\rm TPK}^{\rm enr}\), \(\mathcal F_q\)&
Conservan posición, hoja, orientación, residuo, cociente, acarreo, firma,
frontera, cilindro, prefijo, procedencia y memoria.\\

Selección interna&
\(M_{\rm ph}:\{1,\ldots,9\}\to\{-1,0,+1\}\)&
Activa la evaluación aditiva, multiplicativa o el paso de umbral; no mueve
el estado.\\

Transporte espacial&
\(T_N,T_E,T_S,T_O,F_{\rm loc},F_{\rm obs}\)&
Realiza traslaciones cardinales, realimentación de hoja y cronología de dos
cursores.\\

Lectores modulares&
\(\rho_9,q_9,\rho_3^{\rm or},c_3,q_{1000},r_{1000}\)&
Publican fase, orientación y bloque decimal conservando los cocientes de
levantamiento.\\

Bloque y memoria&
\(\mathcal K_{27}=F_{\rm obs}^{27}\), \(\mathcal N_{27}\)&
El primero compone veintisiete transportes; el segundo filtra eventos de
memoria. Igual índice no implica igual operador.\\

Estado dodecafásico&
\(\mathcal R_{12},C_{12},\Theta_{108},\Gamma_{108}\)&
Distingue registro enriquecido, sector, ley cíclica y soporte ordenado.\\

Periodicidad y profundidades&
\(C_{12}\hookrightarrow C_{108}\twoheadrightarrow C_9\),
\(w_{108},w_{1080}\)&
Conserva la extensión no escindida, las palabras de distintas longitudes y
la memoria vertical.\\

Canales \(90/120\)&
\(\mathcal F_6,\mathcal F_8,\mathcal I_{6\to8},
S_{90},S_{120}\)&
Separa incidencias finitas, colas analíticas y sus normalizaciones.\\

Bidireccionalidad&
\(P_\pm,\operatorname{rev},\operatorname{Ext}_\pm,
N_\pm,\beta,C_M\)&
Realiza intercambio, inversión de ruta, prolongación futura y pasada,
valencia, balance y conjugación.\\

Emisión y coinducción&
\(L_0,L_1,R_{36},G_9,\varprojlim\operatorname{Cyl}_m\)&
Eleva la semilla, abre la frontera y organiza prolongaciones finitas y
compatibles.\\

Lectores numéricos&
\(\mathscr C_\pi,\mathscr P_e,F_{\rm av}\)&
Construyen horquillas racionales, propagación factorial y autoescala de
Fibonacci.\\

Proyección excepcional&
\(\Gamma_{A_W}\), peso, soporte proyectivo&
Transporta la misma emisión a incidencia ternaria sin seleccionar cifras
arquimedianas.\\

Atlas y realización&
firma de ruta, atlas \(81/56/13/324\), \(C_M\)&
Clasifica historias y caracteres antes de aplicar una carta dimensional.\\

Validación finita&
censos, pruebas de necesidad estructural, identidades de integridad y
controles de independencia&
Comprueba cobertura en la profundidad ejecutada y separa construcción,
comprobación independiente y datos de contraste.\\
\bottomrule
\end{longtable}
\endgroup

Las ocho clases describen naturaleza matemática; las catorce agrupaciones,
función de ejecución. Son dos gradaciones del mismo inventario.

\Needspace{.92\textheight}
\subsection{Dominio y tipado de las \(34\) entradas del operador de actualización}

\begin{equation}
 \mathfrak O_{\rm TPK}:=(\mathfrak o_1,\ldots,\mathfrak o_{34}).
 \label{p12-83-c1-s6:eq:u006f-registro-operadores}
\end{equation}

\begingroup
\scriptsize
\renewcommand{\arraystretch}{1.12}
\begin{longtable}{@{}r p{.19\textwidth}p{.28\textwidth}p{.42\textwidth}@{}}
\toprule
\#&Símbolo&Dominio y codominio&Regla o función\\
\midrule
\endfirsthead
\toprule
\#&Símbolo&Dominio y codominio&Regla o función\\
\midrule
\endhead
1&\(\mathcal A_9\)&\((\mathbb Z/9)^2\), matriz y ley interna&
Célula aritmética aditivo--multiplicativa.\\
2&\((\rho_9,q_9)\)&\(\mathbb Z_{>0}\to
\{1,\ldots,9\}\times\mathbb Z_{\geq0}\)&
División nonádica con representante positivo y cociente.\\
3&\(\rho_3^{\rm or},c_3\)&\(\mathbb Z\to
\{-1,0,+1\}\times\mathbb Z\)&
Levantamiento ternario orientado \(n=\rho_3^{\rm or}(n)+3c_3(n)\).\\
4&\(\iota_{1000}:=(q_{1000},r_{1000})\)&\(\mathbb Z\to
\mathbb Z\times\{0,\ldots,999\}\)&
División euclídea \(N=1000q_{1000}+r_{1000}\).\\
5&\(T_d\)&\(X_9\to X_9\), \(d\in\{N,E,S,O\}\)&
Traslación cardinal orientada.\\
6&\(F_{\rm loc}\)&\(\mathcal Q_{\rm loc}\times
\mathscr D_{\rm card}\to\mathcal Q_{\rm loc}\)&
Realimentación local de posición y hoja.\\
7&\(T_{\min}\)&\(\Omega_{\rm seed}\to\mathcal U_6\)&
Emisor mínimo de los dos cursores, incluida la firma coordenada.\\
8&\(G\)&\(\operatorname{im}s_\Sigma\times
\operatorname{im}s_\Pi\to\{0,\ldots,999\}^6\)&
Acoplamiento decimal coordenada a coordenada.\\
9&\(\Psi_6\)&\(\mathcal U_6\to\mathbb F_3^6\)&
Reducción ternaria del catálogo.\\
10&\(\mathcal R_3\)&\(\mathcal Q_{\rm loc}\times
\mathscr D_{\rm card}^3\to\{-1,0,+1\}\)&
Factor observable local de tres pasos.\\
11&\(\mathcal K_{27}\)&\(\mathcal Q_{\rm obs}\to\mathcal Q_{\rm obs}\)&
Composición \(F_{\rm obs}^{27}\).\\
12&\(\mathcal X_{\rm TPK}^{\rm enr}\)&categoría de estados enriquecidos&
Conserva todas las coordenadas tipadas del núcleo.\\
13&\(\mathcal H\)&relación sobre estados enriquecidos&
Transición Hensel que retiene residuo y cociente.\\
14&\(\mathcal C_{3,1000}\)&estados cruzados
\((N,L,D,K,A,B)\)&
Actualización exacta de cilindros ternario--decimales.\\
15&\(\Sigma_{\rm B}\)&memorias de frontera \(\to\) firma conjunta&
Publica la orientación de bloque conservando su registro genealógico.\\
16&\({\rm EF}\text{--}G_9\)&relación de extensión&
Filtra extensiones compatibles y organiza la genealogía nonádica.\\
17&\(\Gamma_9\)&historias enriquecidas \(\longrightarrow\) historias&
Holonomía de una vuelta de fase con memoria de frontera.\\
18&\(X_\chi\)&\(\varprojlim_m\operatorname{Surv}^{(\chi)}_m\)&
Sección global cuando no vaciedad, compatibilidad y unicidad están
demostradas en toda profundidad.\\
19&\(\mathcal P_{\rm APP}^{(2)}\)&categoría bivariada de caminos&
Transporta simultáneamente la evaluación multiplicativa y la acumulación
aditiva.\\
20&\(F_{\rm obs}\)&\(\mathcal Q_{\rm obs}\to\mathcal Q_{\rm obs}\)&
Cronología observable determinista de dos cursores.\\
21&\(\mathcal A_{42}\)&alfabeto de cuarenta y dos tríadas&
Soporte decimal ampliado de la emisión finita.\\
22&\(\operatorname{Ext}_r\)&
\(\operatorname{Ext}_r\subseteq\mathcal F_q\times\mathcal F_{q'}\)&
Relación tipada de prolongación de fibras.\\
23&\(\mathcal K_{\rm ph}\)&\(\mathcal U_6\times C_9\) sobre sí mismo&
Transferencia de continuaciones compatibles con fase explícita.\\
24&\((P,H,Q)\)&\(\mathcal X_\times\to\mathcal X_\times\)&
Avance, media vuelta y cuarto de giro de semillas.\\
25&\((c_{\rm mode},c_{\rm or})\)&caminos \(\to\mathbb Z^2\)&
Cociclos de cambio de modo y orientación.\\
26&\(E_{20}\)&estados finitos \(\to\) veinte bloques&
Prolongación triádica finita con libro de procedencia.\\
27&\(\mathcal R_{12}\)&historia enriquecida \(\to\) registro de doce
sectores&
Sistema temporal orientado con ruta, firma, acarreo y registro genealógico.\\
28&\(R_{\rm sgn}\)&estado terminal \(\to U_{12}^{\rm sgn}
\subseteq\mathbb Z^{12}\)&
Lector del observable armónico firmado.\\
29&\(\mathcal H_{12}\)&\(\mathbb Z^{12}\to\mathbb Z^{12}\)&
Transformación por tres órbitas de paso tres; inversa
\(\mathcal H_{12}/4\) sobre la imagen pertinente.\\
30&\(\mathcal E_{\rm dod}=(D_3,D_4,Q)\)&
\(\mathbb Z^{12}\to\mathbb Z^{25}\)&
Extractor transversal conjunto de aridades tres y cuatro.\\
31&\(\mathcal C_\alpha\)&bloques y acarreos \(\to\) bloques&
Recurrencia dodecafásica de cierre en base mil.\\
32&\((L_{\rm tick},\chi_{\rm mass})\)&historia enriquecida \(\to\)
carácter temporal&
Interfaz dimensional posterior; no participa en el constructor numérico de
esta obra.\\
33&\(\mathcal W_\pm\)&retícula bidireccional de índice doce&
Lectura en ambos sentidos conservando orientación y procedencia.\\
34&\(\mathcal A_{\rm species}\)&familias de extensión \(\to\) atlas&
Interfaz de realización posterior, no un cuarto primitivo.\\
\bottomrule
\end{longtable}
\endgroup

\subsection{Composición causal de la actualización del estado enriquecido}

\begin{definition}[Estado canónico completo]
\label{p12-83-c1-s6:def:u006f-esquema-estado}
Sea
\[
 \mathcal E:=
 \mathbb Z_{\rm quo}\times\mathbb Z_{\rm car}\times
 \mathsf{Mem}\times\mathsf{Term}\times\mathsf{Cyl}\times
 \mathsf{Fr}\times\mathsf{Pref}\times\Lambda .
\]
El estado enriquecido canónico en el instante \(t\) es
\begin{equation}
 x_t=(p_t,m_t,\sigma_t,b_t,r_t;
 q_t,c_t,\eta_t,s_t,C_t,\partial_t,N_t,\lambda_t)
 \in
 X_9\times\{\Sigma,\Pi\}\times\mathsf{TRIT}\times
 C_{12}\times C_9\times\mathcal E.
 \label{p12-83-c1-s6:eq:u006f-esquema-estado}
\end{equation}
Las abreviaturas de estado utilizadas en otras unidades son proyecciones de
esta tupla. En particular, no se omiten el cociente \(q_t\), el acarreo
\(c_t\), el terminal \(s_t\), el cilindro \(C_t\), la frontera
\(\partial_t\), el prefijo \(N_t\) ni el libro \(\lambda_t\).
\end{definition}

Sea \(\rho_t\) la flecha observable determinada por \(x_t\). Las leyes de
transporte de U005 y la transición \(\mathcal C_{3,1000}\) determinan un
endomorfismo de memoria
\begin{align}
 \mathbf E_t:\mathcal E&\longrightarrow\mathcal E,\notag\\
 (q,c,\eta,s,C,\partial,N,\lambda)&\longmapsto
 \left(
 \begin{aligned}
 &\mathsf H(\rho_t)q,
  \mathsf C(\rho_t)c+\kappa(\rho_t),
  \mathsf M(\rho_t)\eta,\\
 &\mathsf T(\rho_t)s,
  \mathsf{Cyl}(\rho_t)C,
  \mathsf{Fr}(\rho_t)\partial,\\
 &\mathsf{Pref}(\rho_t)N,
  \Lambda(\rho_t)\lambda
 \end{aligned}
 \right).
 \label{p12-83-c1-s6:eq:u006f-actualizacion-memoria}
\end{align}
Cada componente publicada en \eqref{p12-83-c1-s6:eq:u006f-actualizacion-memoria} posee el
dominio indicado por su fibra; \(\kappa(\rho_t)\) es la cocadena de acarreo.

Definimos tres endomorfismos del producto canónico. Si
\(\sigma'=M_{\rm ph}(g(t))\) y \(m'=\mu(\sigma',m)\), entonces
\begin{align}
 \mathsf{Sel}_t(p,m,\sigma,b,r;e)
  &:=(p,m',\sigma',b,r;e),\label{p12-83-c1-s6:eq:u006f-selector-tipado}\\
 \mathsf{Tra}_t(p,m,\sigma,b,r;e)
  &:=(T_{d(t)}p,m,\sigma,b,r;e),\label{p12-83-c1-s6:eq:u006f-transporte-tipado}\\
 \mathsf{Upd}_t(p,m,\sigma,b,r;e)
  &:=(p,m,\sigma,b',r';\mathbf E_t(e)),
 \label{p12-83-c1-s6:eq:u006f-memoria-tipada}
\end{align}
donde \((b',r')\) es el destino de \(\rho_t:(b,r)\to(b',r')\). La actualización
causal es la composición de mapas, no la composición de un mapa con un valor
trítico:
\begin{equation}
 \boxed{\mathcal U_t:=
 \mathsf{Upd}_t\circ\mathsf{Tra}_t\circ\mathsf{Sel}_t.}
 \label{p12-83-c1-s6:eq:u006f-tuberia}
\end{equation}

\begin{theorem}[Cierre de la actualización causal]
\label{p12-83-c1-s6:thm:u006f-cierre-tuberia}
La aplicación \(\mathcal U_t\) es un endomorfismo del estado canónico
completo. Los lectores, prolongaciones y proyecciones se aplican después de
\(\mathcal U_t\) sólo cuando su dominio ha sido producido.
\end{theorem}

\begin{proof}
\(\mathsf{Sel}_t\) cambia únicamente modo y TRIT;
\(\mathsf{Tra}_t\) cambia únicamente la posición APP; y
\(\mathsf{Upd}_t\) aplica el producto tipado
\eqref{p12-83-c1-s6:eq:u006f-actualizacion-memoria} y la flecha de fase. Los tres
codominios son el mismo producto de la
\cref{p12-83-c1-s6:def:u006f-esquema-estado}; por tanto su composición está definida y
conserva todas las coordenadas del estado.
\end{proof}

\subsection{Separación tipada entre operadores de transporte, memoria, calendario, extracción y retorno}

El inventario impone
\begin{equation}
 \mathcal S_4\neq D_4\neq\operatorname{sd}_4,
 \qquad
 \mathcal R_{12}\neq C_{12}\neq\Theta_{108}\neq\Gamma_{108},
 \qquad
 F_{\rm obs}\neq\mathcal K_{\rm ph},
 \label{p12-83-c1-s6:eq:u006f-no-identificacion-uno}
\end{equation}
y
\begin{equation}
 U_{12}^{\rm cat}\neq U_{12}^{\rm sgn},
 \qquad
 \mathcal E_{\rm dod}\neq D_{\rm raw}^{90,120}
 \neq\mathcal K_{6\to8}.
 \label{p12-83-c1-s6:eq:u006f-no-identificacion-dos}
\end{equation}
Son incompatibilidades de tipo, no meras desigualdades de valor.

\begin{criterion}[Auditoría de completitud operatoria]
\label{p12-83-c1-s6:crit:u006f-completitud}
Una ejecución es operatoriamente completa sólo si:
\begin{enumerate}
 \item declara las entradas de \(\mathfrak O_{\rm TPK}\) empleadas;
 \item registra el dominio y el codominio de cada composición;
 \item conserva el orden causal de \eqref{p12-83-c1-s6:eq:u006f-tuberia};
 \item distingue dinámica, transferencia, lectura, prolongación y salida;
 \item permite reconstruir desde \(\lambda_t\) cada cambio de bloque y
 cada acarreo;
 \item no promueve una comprobación finita a una sección infinita sin
 probar no vaciedad, compatibilidad y unicidad a toda profundidad.
\end{enumerate}
Cualquier símbolo usado sin tipo o cualquier fusión prohibida por
\eqref{p12-83-c1-s6:eq:u006f-no-identificacion-uno} y
\eqref{p12-83-c1-s6:eq:u006f-no-identificacion-dos} constituye un fallo verificable.
\end{criterion}


% Unidad U006E. Separación tipada de los canales 90/120 y de los símbolos
% epsilon. La realización de incidencia se construye después de W12 y W24.

\section{Descomposición operatorial de los canales de incidencia \texorpdfstring{\(90/120\)}{90/120}: selector temporal, generador nilpotente y componente angular}
\label{p12-83-c1-s7:chap:canales-90-120-epsilon}

Los enteros \(90\) y \(120\) aparecen en varios estratos de la construcción. Su
coincidencia numérica no establece una identidad de objetos. Del mismo modo,
la letra epsilon designa tres construcciones con dominios incompatibles. La
finalidad de esta unidad es fijar sus tipos antes de emplearlos y señalar
con precisión qué realizaciones requieren estructuras que todavía no se han
construido.

\subsection{Operadores temporal, nilpotente y angular: dominios y codominios}

\subsubsection{Acción del selector temporal trítico sobre \(\{-1,0,+1\}\): apertura, umbral y retorno con memoria}

El selector interno del calendario es
\begin{equation}
 \varepsilon_9:\{1,\ldots,9\}
 \longrightarrow\{-1,0,+1\},
 \qquad
 \varepsilon_9(r)=
 \begin{cases}
  +1,&r\in\{1,4,7\},\\
  -1,&r\in\{2,5,8\},\\
  0,&r\in\{3,6,9\}.
 \end{cases}
 \label{p12-83-c1-s7:eq:u006e-epsilon-temporal}
\end{equation}
Su codominio es el TRIT orientado. Decide qué hoja actúa y no desplaza por
sí mismo ninguna posición.

\subsubsection{Generador parabólico nilpotente \(J_0^2=0\) y evolución en el umbral temporal}

La clase
\begin{equation}
 \varepsilon_{\rm nil}
 \in
 \mathcal Q_{\rm n}:=
 \mathbb F_3[\varepsilon_{\rm nil}]
 /(\varepsilon_{\rm nil}^{\,2})
 \label{p12-83-c1-s7:eq:u006e-epsilon-nilpotente}
\end{equation}
satisface
\begin{equation}
 \varepsilon_{\rm nil}\neq0,
 \qquad
 \varepsilon_{\rm nil}^{\,2}=0.
 \label{p12-83-c1-s7:eq:u006e-epsilon-nilpotente-relacion}
\end{equation}
Es un elemento de un álgebra cuadrática degenerada; no es un selector
temporal.

\subsubsection{Residuo angular inducido después de la selección temporal y conservación de orientación}

El símbolo
\begin{equation}
 \varepsilon_{\rm res}^{\circ}\in\mathbb R
 \label{p12-83-c1-s7:eq:u006e-epsilon-residual-tipo}
\end{equation}
designa, cuando se introduzcan sus parámetros geométricos, el residuo de una
identidad angular. En esta etapa sólo queda reservado su tipo. No se lo
obtiene desde \(\varepsilon_9\) ni desde
\(\varepsilon_{\rm nil}\), y no interviene en la construcción de
\(\pi,e,\varphi\) o \(\alpha\).

\begin{proposition}[No identificación de los símbolos epsilon]
\label{p12-83-c1-s7:prop:u006e-no-identificacion-epsilon}
No existe una igualdad bien tipada entre
\[
 \varepsilon_9,\qquad
 \varepsilon_{\rm nil},\qquad
 \varepsilon_{\rm res}^{\circ}.
\]
\end{proposition}

\begin{proof}
El primer objeto es una función entre conjuntos finitos; el segundo, un
elemento nilpotente de un álgebra sobre \(\mathbb F_3\); el tercero, un
escalar real con unidad angular. Sus dominios, codominios y leyes de
composición son distintos.
\end{proof}

\subsection{Distribución de los índices \(90/120\) entre los \(6\) objetos del sistema de incidencia}

\begin{definition}[Extractor dodecafásico]
\label{p12-83-c1-s7:def:u006e-extractor-dodecafasico}
El extractor integral
\begin{equation}
 \mathcal E_{\rm dod}:=(D_3,D_4,Q):
 \mathbb Z^{12}\longrightarrow\mathbb Z^{25}
 \label{p12-83-c1-s7:eq:u006e-extractor-dodecafasico}
\end{equation}
lee diferencias de aridades tres y cuatro y una coordenada de suma. Los
subíndices no son longitudes angulares.
\end{definition}

\begin{definition}[Diferencia cruda de series]
\label{p12-83-c1-s7:def:u006e-diferencia-cruda}
Para \(|q|<1\), se define
\begin{equation}
 S_m^{\rm raw}(q):=\frac{q^m}{1-q^{3m}},
 \qquad
 D_{\rm raw}^{90,120}(q)
 :=S_{90}^{\rm raw}(q)-S_{120}^{\rm raw}(q).
 \label{p12-83-c1-s7:eq:u006e-series-crudas}
\end{equation}
Este objeto es una función analítica de \(q\); no es un censo finito.
\end{definition}

\begin{definition}[Familias marcadas de incidencia]
\label{p12-83-c1-s7:def:u006e-familias-incidencia}
Una vez construidas una hexada \(H\), su imagen
\(\ell(H)\) en una dodecada binaria \(D\), y la octada elevada
\(O_H\), se definirán
\begin{align}
 \mathcal F_6(H)
 &:=\ell(H)\times\binom{\ell(H)}2,
 \label{p12-83-c1-s7:eq:u006e-familia-seis}\\
 \mathcal F_8(H)
 &:=O_H\times\binom{\ell(H)}2.
 \label{p12-83-c1-s7:eq:u006e-familia-ocho}
\end{align}
Por construcción,
\begin{equation}
 |\mathcal F_6(H)|
 =6\binom62=90,
 \qquad
 |\mathcal F_8(H)|
 =8\binom62=120.
 \label{p12-83-c1-s7:eq:u006e-cardinales-incidencia}
\end{equation}
La inclusión \(\ell(H)\hookrightarrow O_H\) inducirá
\begin{equation}
 \mathcal I_{6\to8}:
 \mathcal F_6(H)\hookrightarrow\mathcal F_8(H).
 \label{p12-83-c1-s7:eq:u006e-inclusion-incidencia}
\end{equation}
\end{definition}

Las \cref{p12-83-c1-s7:eq:u006e-familia-seis,p12-83-c1-s7:eq:u006e-familia-ocho} fijan desde ahora el
tipo. Su existencia efectiva se demostrará después de construir
\(C_W\), \(W_{12}\), el dato binario \(G_{24},D,\ell\) y \(W_{24}\).
No se anticipa aquí esa rama.

\begin{definition}[Combinación normalizada de colas]
\label{p12-83-c1-s7:def:u006e-combinacion-colas}
Para parámetros \(q_-,q_+\in(0,1)\), sean
\begin{equation}
 S_m(q):=\frac{q^m}{1-q^{3m}},
 \qquad
 \Delta S_m:=S_m(q_-)-S_m(q_+).
 \label{p12-83-c1-s7:eq:u006e-colas-normalizadas}
\end{equation}
El funcional
\begin{equation}
 \mathcal K_{6\to8}
 :=12\,\Delta S_{90}-\Delta S_{120}
 \label{p12-83-c1-s7:eq:u006e-funcional-seis-ocho}
\end{equation}
es una combinación analítica de colas. Su coeficiente doce pertenece a su
normalización y no lo convierte en el registro
\(\mathcal R_{12}\).
\end{definition}

\begin{definition}[Semigrupo de grados y momentos]
\label{p12-83-c1-s7:def:u006e-semigrupo-momentos}
El semigrupo aditivo generado por ambos grados es
\begin{equation}
 \mathfrak S_{90,120}
 :=\langle90,120\rangle
 =\{90a+120b:a,b\in\mathbb N_0\}.
 \label{p12-83-c1-s7:eq:u006e-semigrupo}
\end{equation}
Para \(n\in\mathfrak S_{90,120}\), el momento orientado es
\begin{equation}
 s_n:=q_-^n-q_+^n.
 \label{p12-83-c1-s7:eq:u006e-momento}
\end{equation}
\end{definition}

Los momentos son monomios diferenciales indexados por un semigrupo; las
series de \eqref{p12-83-c1-s7:eq:u006e-colas-normalizadas} son sumas infinitas de esos
momentos sobre progresiones específicas.

\subsection{Defecto normalizado de la incidencia}

Cuando la inclusión \(\mathcal I_{6\to8}\) haya sido construida, su
diferencia de cardinalidades será
\[
 |\mathcal F_6(H)|-|\mathcal F_8(H)|=90-120=-30.
\]
Si se declara la normalización
\[
 24\binom62=360,
\]
se obtiene
\begin{equation}
 \frac{90-120}{24\binom62}
 =\frac{6-8}{24}
 =-\frac1{12}.
 \label{p12-83-c1-s7:eq:u006e-defecto-normalizado}
\end{equation}
La inclusión determina \(-30\); el racional \(-1/12\) depende además de
la normalización declarada. No se lo debe presentar como si la cardinalidad
por sí sola fijara el denominador.

\subsection{Invariantes de tipado}

\begin{theorem}[Separación de los seis objetos]
\label{p12-83-c1-s7:thm:u006e-separacion-seis-objetos}
Los objetos
\begin{equation}
 \mathcal E_{\rm dod},\quad
 D_{\rm raw}^{90,120},\quad
 \mathcal I_{6\to8},\quad
 \mathcal K_{6\to8},\quad
 \mathfrak S_{90,120},\quad
 \varepsilon_{\rm res}^{\circ}
 \label{p12-83-c1-s7:eq:u006e-seis-objetos}
\end{equation}
son mutuamente distintos por tipo: son, respectivamente, un lector entero,
una función analítica, un morfismo finito, una combinación de series, un
semigrupo de índices y un escalar real angular.
\end{theorem}

\begin{proof}
La conclusión se obtiene comparando sus dominios y codominios:
\[
 \mathbb Z^{12}\to\mathbb Z^{25},\quad
 \mathbb D\to\mathbb C,\quad
 \mathcal F_6(H)\to\mathcal F_8(H),\quad
 (0,1)^2\to\mathbb R,\quad
 \mathbb N_0^2\to\mathbb N_0,\quad
 \{\ast\}\to\mathbb R.
\]
No hay dos firmas iguales.
\end{proof}

\begin{criterion}[Criterios de refutación]
\label{p12-83-c1-s7:crit:u006e-refutacion}
La clasificación queda refutada si se altera alguno de sus dominios o
codominios, si el censo de
\eqref{p12-83-c1-s7:eq:u006e-cardinales-incidencia} no reproduce \(90\) y \(120\), si
se atribuye \(-1/12\) a la inclusión sin declarar la normalización, o si
se presenta el residual angular como salida del generador de constantes de
esta monografía.
\end{criterion}


% Unidad U006A. Factor local, observacion estroboscopica y transicion exacta
% de cilindros.

\section{Factor local de transporte y transición ternario--decimal}
\label{p12-83-c1-s8:chap:tpk-factor-local-transicion-cilindros}

Esta unidad separa tres operaciones que intervienen en el transporte del
Topological Phase Kernel y que no comparten dominio: la realización cardinal
de un trit, el muestreo estroboscópico de una sucesión y la actualización
exacta de un par de cilindros ternario--decimales. La primera es una dinámica
finita; la segunda es un lector temporal; la tercera conserva los enteros
que permiten decidir compatibilidad sin aritmética de punto flotante.

\subsection{Estado local y realimentación de hoja}

Sea
\[
 X_9=(\mathbb Z/9\mathbb Z)^2,
 \qquad
 \mathscr D_{\rm card}=\{N,E,S,O\},
\]
con
\[
 \delta(N)=(-1,0),\quad \delta(E)=(0,1),\quad
 \delta(S)=(1,0),\quad \delta(O)=(0,-1).
\]
Las evaluaciones visibles de APP son
\[
 L_\Sigma(x,y)=\rho_9((x+1)+(y+1)),
 \qquad
 L_\Pi(x,y)=\rho_9((x+1)(y+1)),
\]
donde \(\rho_9(n)=1+((n-1)\bmod9)\). Definimos además
\[
 \rho_3^{\rm or}([0])=0,
 \qquad
 \rho_3^{\rm or}([1])=+1,
 \qquad
 \rho_3^{\rm or}([2])=-1.
\]

\begin{definition}[Factor local de APP--TRIT]
\label{p12-83-c1-s8:def:u006a-factor-local}
El espacio de estados locales es
\[
 \mathcal Q_{\rm loc}=X_9\times\{\Sigma,\Pi\},
 \qquad |\mathcal Q_{\rm loc}|=162.
\]
Para \(q=(p,m)\), \(d\in\mathscr D_{\rm card}\), ponemos
\[
 p_d=p+\delta(d),
 \qquad
 a_d(q)=\rho_3^{\rm or}([L_m(p_d)]_3),
\]
y
\[
 \mu(a,m)=
 \begin{cases}
  \Sigma,&a=+1,\\
  \Pi,&a=-1,\\
  m,&a=0.
 \end{cases}
\]
La transición local es
\[
 F_d^{\rm loc}(p,m)=(p_d,\mu(a_d(p,m),m)).
\]
\end{definition}

La posición se modifica antes de evaluar la hoja. El trit nulo conserva el
último modo activo; no elimina la hoja ni reinicializa la historia.

Para una palabra cardinal \(u=d_1d_2d_3\), sea
\[
 q_j=F_{d_j}^{\rm loc}(q_{j-1}),\qquad q_0=q,
\]
y denotemos por
\[
 R_3(q,u)\in\{-1,0,+1\}
\]
el trit producido en el tercer desplazamiento.

\begin{definition}[Observación de paso tres]
\label{p12-83-c1-s8:def:u006a-obs3}
Sobre una sucesión orientada \(a=(a_t)_{t\geq1}\) definimos
\[
 \operatorname{Obs}_3(a)=(a_3,a_6,a_9,\ldots).
\]
El mapa \(R_3\) toma un estado y una palabra cardinal de longitud tres;
\(\operatorname{Obs}_3\) toma una sucesión y publica una subsucesión. Son
operadores de tipos distintos.
\end{definition}

\begin{theorem}[Sobreyectividad exacta del factor muestreado]
\label{p12-83-c1-s8:thm:u006a-sobreyectividad-r3}
Para todo \(q\in\mathcal Q_{\rm loc}\) y todo
\(b\in\{-1,0,+1\}\), la fibra
\[
 \mathcal R_3(q,b)
 =\{u\in\mathscr D_{\rm card}^3:R_3(q,u)=b\}
\]
es no vacía y satisface
\[
 \boxed{8\leq |\mathcal R_3(q,b)|\leq40.}
\]
En particular, las \(162\cdot3=486\) fibras estado--salida son no vacías.
\end{theorem}

\begin{proof}
Desde cada uno de los \(162\) estados se enumeran las \(4^3=64\) palabras
cardinales y se aplica literalmente la recurrencia de la
\cref{p12-83-c1-s8:def:u006a-factor-local}. El censo exhaustivo de los \(486\) pares
\((q,b)\) tiene mínimo ocho, máximo cuarenta y ningún valor cero. El
enunciado es, por tanto, una identidad finita sobre el sistema definido, no
una hipótesis de mezcla.
\end{proof}

\begin{corollary}[Realización de una sucesión orientada arbitraria]
\label{p12-83-c1-s8:cor:u006a-realizacion-sucesion}
Para toda sucesión
\((b_k)_{k\geq1}\in\{-1,0,+1\}^{\mathbb N}\) existe una historia cardinal
cuya observación satisface
\[
 (a_3,a_6,a_9,\ldots)=(b_1,b_2,b_3,\ldots).
\]
\end{corollary}

\begin{proof}
Fijado un orden total de \(\mathscr D_{\rm card}^3\), en la etapa \(k\) se
toma el primer elemento de \(\mathcal R_3(q_{k-1},b_k)\). El estado terminal
del bloque vuelve a pertenecer a \(\mathcal Q_{\rm loc}\), de modo que la
elección puede iterarse indefinidamente.
\end{proof}

El corolario demuestra capacidad de realización del factor muestreado. No
selecciona las secciones de \(\pi,e,\varphi\), no fija los dos trits
intermedios y no sustituye la compatibilidad del estado enriquecido.

\subsection{Cilindros de los \(2\) lectores}

Una palabra ternaria de longitud \(L\), interpretada como el entero \(N\),
determina
\[
 I_3(N,L)=\left[\frac{N}{3^L},\frac{N+1}{3^L}\right).
\]
Un prefijo de \(K\) tríadas decimales, interpretado como el entero \(D\),
determina
\[
 I_{1000}(D,K)=
 \left[\frac{D}{1000^K},\frac{D+1}{1000^K}\right).
\]
Introducimos los residuos cruzados
\[
 A=N1000^K-D3^L,
 \qquad
 B=(N+1)1000^K-D3^L=A+1000^K.
\]

\begin{lemma}[Criterio entero de compatibilidad]
\label{p12-83-c1-s8:lem:u006a-criterio-entero}
Los cilindros \(I_3(N,L)\) e \(I_{1000}(D,K)\) se intersectan si y sólo si
\[
 \boxed{A<3^L,\qquad B>0.}
\]
Equivalentemente,
\[
 N1000^K<(D+1)3^L,
 \qquad
 (N+1)1000^K>D3^L.
\]
\end{lemma}

\begin{proof}
Dos intervalos semiabiertos \([a,b)\) y \([c,d)\) se intersectan
exactamente cuando \(a<d\) y \(c<b\). La multiplicación por
\(3^L1000^K>0\) produce las dos desigualdades. Restar \(D3^L\) y
\(N1000^K\), respectivamente, da la forma en \(A,B\).
\end{proof}

\subsection{Operador de completación \(C_{3,1000}\): extensión ternaria y conservación de la unidad en \(10^3\) posiciones}

Un nuevo bloque de seis trits se representa por
\[
 u\in\{0,\ldots,3^6-1\}=\{0,\ldots,728\}.
\]
La actualización ternaria es
\[
 N'=729N+u,
 \qquad L'=L+6.
\]
La publicación decimal tiene dos casos. Si el cilindro refinado todavía no
determina una nueva tríada, entonces \(K'=K\), \(D'=D\). Si determina
\(\delta\in\{0,\ldots,999\}\), entonces
\[
 K'=K+1,
 \qquad D'=1000D+\delta.
\]

\begin{definition}[Estado cruzado y transición exacta]
\label{p12-83-c1-s8:def:u006a-c31000}
El estado de cilindros es la séxtupla
\[
 \mathfrak c=(N,L,D,K,A,B)
\]
sujeta a
\[
 A=N1000^K-D3^L,
 \qquad B=A+1000^K.
\]
La transición \(\mathcal C_{3,1000}\) adjunta \(u\) y, cuando corresponde,
\(\delta\), y actualiza
\[
 (N,L,D,K,A,B)\longmapsto
 (N',L',D',K',A',B')
\]
mediante
\[
 A'=
 \begin{cases}
  729A+u1000^K,&K'=K,\\[1mm]
  729000A+u1000^{K+1}-\delta\,729\,3^L,&K'=K+1,
 \end{cases}
 \qquad
 B'=A'+1000^{K'}.
\]
La prolongación es admisible exactamente cuando
\[
 A'<3^{L'},\qquad B'>0.
\]
\end{definition}

\begin{theorem}[Exactitud y suficiencia del estado cruzado]
\label{p12-83-c1-s8:thm:u006a-exactitud-c31000}
La transición de la \cref{p12-83-c1-s8:def:u006a-c31000} conserva exactamente las
identidades
\[
 A'=N'1000^{K'}-D'3^{L'},
 \qquad
 B'=(N'+1)1000^{K'}-D'3^{L'}.
\]
Por consiguiente, \((A',B',L',K')\) basta para decidir la compatibilidad del
nuevo par de cilindros, mientras \((N',D')\) conserva su procedencia
reconstructiva.
\end{theorem}

\begin{proof}
Si no aparece tríada decimal,
\[
 N'1000^K-D3^{L+6}
 =(729N+u)1000^K-729D3^L
 =729A+u1000^K.
\]
Si aparece \(\delta\),
\begin{align*}
 N'1000^{K+1}-D'3^{L+6}
 &=(729N+u)1000^{K+1}-(1000D+\delta)7293^L\\
 &=729000A+u1000^{K+1}-\delta\,729\,3^L.
\end{align*}
En ambos casos, aumentar \(N'\) en una unidad añade \(1000^{K'}\), lo que
da \(B'\). El criterio de la \cref{p12-83-c1-s8:lem:u006a-criterio-entero} concluye.
\end{proof}

\begin{proposition}[Compatibilidad con la concatenación]
\label{p12-83-c1-s8:prop:u006a-concatenacion-c31000}
Aplicar sucesivamente \(\mathcal C_{3,1000}\) a dos bloques produce el
mismo par \((N'',D'')\) y los mismos residuos cruzados que adjuntar de una
vez las dos palabras ternarias y las tríadas decimales publicadas en el
mismo orden.
\end{proposition}

\begin{proof}
Las recurrencias \(N\mapsto729N+u\) y
\(D\mapsto1000D+\delta\) son las leyes de concatenación posicional en
bases \(729\) y \(1000\). El
\cref{p12-83-c1-s8:thm:u006a-exactitud-c31000} identifica \(A,B\) con funciones de esos
cuatro enteros; por tanto, su actualización por etapas coincide con la
evaluación posterior a la concatenación.
\end{proof}

\subsection{Descomposición del Topological Phase Kernel en operadores de avance, giro, acarreo y memoria}

La cadena correcta es
\[
 (q,u)\xmapsto{\ R_3\ }b,
 \qquad
 (a_t)_{t\geq1}
 \xmapsto{\ \operatorname{Obs}_3\ }
 (a_{3k})_{k\geq1},
 \qquad
 (\mathfrak c,u,\delta)
 \xmapsto{\ \mathcal C_{3,1000}\ }
 \mathfrak c'.
\]
\(R_3\) prueba realizabilidad local; \(\operatorname{Obs}_3\) publica una
subsucesión; \(\mathcal C_{3,1000}\) transporta compatibilidad entre dos
bases. Ninguno de los tres puede reemplazar a la relación de supervivencia
del TPK completo.

\begin{criterion}[Criterios de refutación]
\label{p12-83-c1-s8:crit:u006a-refutacion}
La construcción queda refutada en su dominio si ocurre cualquiera de los
hechos siguientes:
\begin{enumerate}
 \item alguna de las \(486\) fibras \(\mathcal R_3(q,b)\) es vacía o queda
       fuera del intervalo \([8,40]\);
 \item el trit nulo cambia el modo activo;
 \item se identifica \(R_3\) con \(\operatorname{Obs}_3\);
 \item la recurrencia de \(A'\) no reproduce su definición directa;
 \item una prolongación declarada compatible viola \(A'<3^{L'}\) o
       \(B'>0\);
 \item la decisión depende de una tolerancia numérica no presente en las
       desigualdades enteras.
\end{enumerate}
\end{criterion}


% Unidad U006D. Registro dodecafásico integral y descomposición 1+5+6.

\section{Registro dodecafásico integral, extensión cíclica y descomposición \texorpdfstring{\(1+5+6\)}{1+5+6}}
\label{p12-83-c1-s9:chap:registro-dodecafasico-integral}

El periodo observable de ciento ocho posiciones contiene doce sectores de
nueve posiciones. Esta igualdad de cardinales no basta para describir el
estado temporal: deben distinguirse el soporte ordenado, la ley cíclica, el
índice sectorial y el registro enriquecido de transportes, orientaciones,
acarreos e incidencias. Esta unidad construye los cuatro objetos y demuestra
la reconstrucción integral de las doce coordenadas de memoria.

\subsection{Soporte ordenado y coordenada cíclica}

Para \(m\in\{1,\ldots,12\}\), definimos
\begin{equation}
 E_m:=\{9(m-1)+1,\ldots,9m\}.
 \label{p12-83-c1-s9:eq:u006d-bloques-nonadicos}
\end{equation}
El soporte ordenado y el grupo cíclico son
\begin{equation}
 \Gamma_{108}:=
 E_1\sqcup\cdots\sqcup E_{12},
 \qquad
 \Theta_{108}:=\mathbb Z/108\mathbb Z.
 \label{p12-83-c1-s9:eq:u006d-gamma-theta}
\end{equation}
\(\Gamma_{108}\) conserva posición, orden y pertenencia a bloque;
\(\Theta_{108}\) conserva la ley de traslación. No son el mismo objeto.

Para \(\theta\in\Theta_{108}\), sea
\(\widetilde\theta\in\{0,\ldots,107\}\) su representante. Definimos
\begin{equation}
 \beta(\theta):=
 \left[
 \left\lfloor\frac{\widetilde\theta}{9}\right\rfloor
 \right]_{12},
 \qquad
 r(\theta):=1+(\widetilde\theta\bmod9).
 \label{p12-83-c1-s9:eq:u006d-coordenadas-bloque-residuo}
\end{equation}
La primera coordenada localiza el sector dodecafásico; la segunda publica
la posición positiva \(1,\ldots,9\) dentro del sector.

\subsection{Extensión cíclica no escindida}

Sean \(C_n=\mathbb Z/n\mathbb Z\). Las aplicaciones
\begin{equation}
 \iota:C_{12}\longrightarrow C_{108},
 \quad \iota([b]_{12})=[9b]_{108},
 \qquad
 p:C_{108}\longrightarrow C_9,
 \quad p([\theta]_{108})=[\theta]_9
 \label{p12-83-c1-s9:eq:u006d-mapas-extension}
\end{equation}
forman la sucesión
\begin{equation}
 0\longrightarrow C_{12}
 \xrightarrow{\;\iota\;}C_{108}
 \xrightarrow{\;p\;}C_9
 \longrightarrow0.
 \label{p12-83-c1-s9:eq:u006d-extension}
\end{equation}
Para la sección conjuntista
\[
 s([r]_9):=[\widetilde r]_{108},
 \qquad 0\leq\widetilde r<9,
\]
el defecto de aditividad es
\begin{equation}
 c(r,r'):=
 \left[
 \left\lfloor
 \frac{\widetilde r+\widetilde r'}9
 \right\rfloor
 \right]_{12},
 \qquad
 s(r)+s(r')
 =s(r+r')+\iota(c(r,r')).
 \label{p12-83-c1-s9:eq:u006d-cociclo-extension}
\end{equation}

\begin{theorem}[Exactitud, no escisión y clase cohomológica]
\label{p12-83-c1-s9:thm:u006d-extension-no-escindida}
La sucesión \eqref{p12-83-c1-s9:eq:u006d-extension} es exacta y no se escinde. Para la
acción trivial,
\begin{equation}
 H^2(C_9,C_{12})
 \simeq C_{12}/9C_{12}
 \simeq C_3,
 \label{p12-83-c1-s9:eq:u006d-cohomologia-extension}
\end{equation}
y el cociclo \eqref{p12-83-c1-s9:eq:u006d-cociclo-extension} representa su clase
generadora.
\end{theorem}

\begin{proof}
La imagen de \(\iota\) es el subgrupo de múltiplos de nueve, que coincide
con \(\ker p\); \(p\) es sobreyectiva. Esto prueba la exactitud. Si la
extensión se escindiera, \(C_{108}\) sería isomorfo a
\(C_{12}\times C_9\), pero
\[
 \exp(C_{108})=108,
 \qquad
 \exp(C_{12}\times C_9)=\operatorname{mcm}(12,9)=36.
\]
La identidad de cociclo es la división euclídea de
\(\widetilde r+\widetilde r'\) por nueve. Para un grupo cíclico con acción
trivial,
\(H^2(C_9,C_{12})\simeq C_{12}/9C_{12}\); el acarreo unitario proyecta a
\([1]\), que genera el cociente de orden tres.
\end{proof}

La no escisión tipa el retorno:
\[
 g_0(\theta+9)=g_0(\theta),
 \qquad
 \beta(\theta+9)=\beta(\theta)+1\pmod{12}.
\]
Nueve iteraciones cierran la fase visible y trasladan un sector; ciento
ocho cierran las dos coordenadas del reloj. El estado enriquecido conserva
además la memoria vertical, el cilindro, la frontera y la genealogía.

\subsection{Registro temporal enriquecido: codificación conjunta de fase, hoja, frontera y memoria}

\begin{definition}[Registro dodecafásico orientado]
\label{p12-83-c1-s9:def:u006d-registro-dodecafasico}
Para una historia enriquecida, sea \(\tau_m\) la subruta ordenada sobre
\(E_m\), \(\sigma_m\) su orientación, \(\kappa_m\) el acarreo que cruza la
frontera del bloque y \(\Lambda_m\) su libro local de incidencias. El
registro completo es
\begin{equation}
 \mathcal R_{12}:=
 \bigl(
 (E_m,\tau_m,\sigma_m,\kappa_m,\Lambda_m)
 \bigr)_{m=1}^{12}.
 \label{p12-83-c1-s9:eq:u006d-registro-r12}
\end{equation}
\end{definition}

Se distinguen así cuatro niveles:
\begin{equation}
 \boxed{
 \mathcal R_{12}
 \neq C_{12}
 \neq\Theta_{108}
 \neq\Gamma_{108}.}
 \label{p12-83-c1-s9:eq:u006d-cuatro-objetos}
\end{equation}
El primero es un registro enriquecido; el segundo, un índice de sector; el
tercero, un grupo de traslaciones; el cuarto, un soporte ordenado. Existen
proyecciones entre ellos, pero no identificaciones.

Sea \(q_{\rm mem}\in\mathbb Z_{\geq0}\) el número no acotado de vueltas
nonádicas. La coordenada dodecafásica es sólo su residuo
\[
 \beta=[q_{\rm mem}]_{12}.
\]
Tras ciento ocho pasos,
\[
 \beta\longmapsto\beta,
 \qquad
 q_{\rm mem}\longmapsto q_{\rm mem}+12.
\]
El calendario retorna; la memoria vertical no se reinicia.

\subsection{Extractor firmado de eventos}

Sea \((d_t)_{t=1}^{108}\) el registro digital sobre
\(\Gamma_{108}\). Definimos
\begin{equation}
 \mathcal D_{\rm evt}:=\{2,4,6,8\},
 \qquad
 \Sigma_{12}:=(+,+,+,-,-,-,+,+,+,-,-,-).
 \label{p12-83-c1-s9:eq:u006d-datos-eventos}
\end{equation}
Para \(m=0,\ldots,11\), sea
\begin{equation}
 e_m:=
 \Sigma_{12}(m+1)\,
 \#\{t\in E_{m+1}:d_t\in\mathcal D_{\rm evt}\}
 \in\mathbb Z.
 \label{p12-83-c1-s9:eq:u006d-eventos-firmados}
\end{equation}
La aplicación conserva el censo local y la orientación sectorial como
factores separados.

Emparejamos las dos semicoronas mediante
\begin{equation}
 p_i:=e_i+e_{i+6},
 \qquad
 a_i:=e_i-e_{i+6},
 \qquad 0\leq i\leq5.
 \label{p12-83-c1-s9:eq:u006d-pares-antipodales}
\end{equation}
Definimos el censo total, las cinco diferencias respecto del sexto par y
las seis antisimetrías:
\begin{equation}
 s_C:=\sum_{i=0}^{5}p_i,
 \qquad
 M_i:=p_i-p_5\quad(0\leq i\leq4),
 \qquad
 A_6:=(a_0,\ldots,a_5).
 \label{p12-83-c1-s9:eq:u006d-coordenadas-uno-cinco-seis}
\end{equation}

\begin{definition}[Lector integral \(1+5+6\)]
\label{p12-83-c1-s9:def:u006d-lector-integral}
El lector lineal es
\begin{equation}
 \mathcal L_{12}:\mathbb Z^{12}\longrightarrow\mathbb Z^{12},
 \qquad
 \mathcal L_{12}(e)
 :=
 (s_C,M_0,\ldots,M_4,a_0,\ldots,a_5).
 \label{p12-83-c1-s9:eq:u006d-lector-l12}
\end{equation}
La distribución de coordenadas es
\begin{equation}
 \underbrace{1}_{s_C}
 +\underbrace{5}_{M_0,\ldots,M_4}
 +\underbrace{6}_{a_0,\ldots,a_5}
 =12.
 \label{p12-83-c1-s9:eq:u006d-uno-cinco-seis}
\end{equation}
\end{definition}

\begin{theorem}[Reconstrucción exacta e imagen integral]
\label{p12-83-c1-s9:thm:u006d-inversa-l12}
El lector \(\mathcal L_{12}\) es inyectivo. Sobre su imagen,
\begin{align}
 p_5&=
 \frac{s_C-\sum_{i=0}^{4}M_i}{6},
 &
 p_i&=p_5+M_i\quad(0\leq i\leq4),
 \label{p12-83-c1-s9:eq:u006d-inversa-p}\\
 e_i&=\frac{p_i+a_i}{2},
 &
 e_{i+6}&=\frac{p_i-a_i}{2}
 \quad(0\leq i\leq5).
 \label{p12-83-c1-s9:eq:u006d-inversa-e}
\end{align}
Una tupla de salida pertenece a \(\operatorname{im}\mathcal L_{12}\) si y
sólo si
\begin{equation}
 s_C-\sum_{i=0}^{4}M_i\equiv0\pmod6
 \label{p12-83-c1-s9:eq:u006d-condicion-divisibilidad}
\end{equation}
y, para los \(p_i\) reconstruidos,
\begin{equation}
 p_i\equiv a_i\pmod2
 \qquad(0\leq i\leq5).
 \label{p12-83-c1-s9:eq:u006d-condicion-paridad}
\end{equation}
\end{theorem}

\begin{proof}
De \(M_i=p_i-p_5\) se obtiene
\[
 s_C=\sum_{i=0}^{4}(p_5+M_i)+p_5
 =6p_5+\sum_{i=0}^{4}M_i,
\]
que da la primera ecuación de \eqref{p12-83-c1-s9:eq:u006d-inversa-p}; las demás son
inmediatas. Sumar y restar
\(p_i=e_i+e_{i+6}\) y \(a_i=e_i-e_{i+6}\) produce
\eqref{p12-83-c1-s9:eq:u006d-inversa-e}. La divisibilidad por seis y las seis
congruencias de paridad son precisamente las condiciones necesarias y
suficientes para que todas las coordenadas reconstruidas sean enteras.
\end{proof}

La división por seis aparece después del censo orientado, no durante él. Las
condiciones integrales de esta unidad son exactamente
\eqref{p12-83-c1-s9:eq:u006d-condicion-divisibilidad} y
\eqref{p12-83-c1-s9:eq:u006d-condicion-paridad}; no se atribuye aquí una forma normal de
Smith a una matriz que no haya sido previamente definida.

\subsection{\(2\) palabras dodecafásicas de distinta procedencia}

Debe conservarse la distinción
\begin{equation}
 U_{12}^{\rm cat}:=U_6\Vert U_6,
 \qquad
 U_{12}^{\rm sgn}:=\mathcal H_{12}(K).
 \label{p12-83-c1-s9:eq:u006d-dos-palabras-doce}
\end{equation}
La primera concatena dos emisiones catalogales de longitud seis. La segunda
es una coordenada firmada ligada reversiblemente a \(K\) por el operador
\(\mathcal H_{12}\). Su longitud común no implica una misma procedencia,
semántica o ley de transporte.

\begin{criterion}[Criterios de refutación]
\label{p12-83-c1-s9:crit:u006d-refutacion}
La construcción queda refutada si la sucesión
\eqref{p12-83-c1-s9:eq:u006d-extension} no es exacta, si su clase cohomológica es
trivial, si una tupla de la imagen viola
\eqref{p12-83-c1-s9:eq:u006d-condicion-divisibilidad} o
\eqref{p12-83-c1-s9:eq:u006d-condicion-paridad}, si las fórmulas inversas no recuperan
los doce enteros iniciales, o si el retorno de \(\beta\) se interpreta como
reinicio de \(q_{\rm mem}\).
\end{criterion}


% Unidad U006B. Funtor ciclotomico y elevacion modular de caminos.

\section{Funtor ciclotómico de fase--longitud y elevación modular de caminos}
\label{p12-83-c1-s10:chap:funtor-ciclotomico-elevacion-caminos}

El cociente odométrico \(\mathcal O_{9,12}\) conserva simultáneamente la
posición nonádica, el registro dodecafásico y la longitud entera de cada
flecha. En esta unidad se realiza su fase sobre fibras \(A_2\) y se prueba
la elevación cardinal que sustituye cada arista por cuatro aristas
orientadas. La representación ciclotómica y la longitud absoluta se
mantendrán como coordenadas distintas.

\subsection{Representación de Coxeter de orden \(3\) sobre las fibras dodecafásicas de tipo \(A_2\)}

Sea
\begin{equation}
 G_{A_2}:=
 \begin{pmatrix}2&-1\\-1&2\end{pmatrix},
 \qquad
 C:=
 \begin{pmatrix}0&-1\\1&-1\end{pmatrix}.
 \label{p12-83-c1-s10:eq:u006b-dato-coxeter}
\end{equation}
Por el \cref{p12-83-c1-s5:prop:coxeter-a2-autonomo},
\begin{equation}
 C^3=I_2,
 \qquad
 C^{\mathsf T}G_{A_2}C=G_{A_2},
 \qquad
 \det(XI_2-C)=X^2+X+1.
 \label{p12-83-c1-s10:eq:u006b-identidades-coxeter}
\end{equation}
Así, \(C\) realiza el carácter ciclotómico de \(C_3\) sobre \(A_2\).

Para cada \(b\in\mathbb Z/12\mathbb Z\), sea \(A_{2,b}\) una copia de
\(A_2\), sea
\[
 \iota_b:A_{2,b}\xrightarrow{\;\sim\;}A_2
\]
una identificación integral y sea
\(P_b\in O(G_{A_2})\) un marco que depende de \(b\), pero no del residuo
\(r\). Definimos
\begin{equation}
 \rho_b:=
 \iota_b^{-1}P_bCP_b^{-1}\iota_b,
 \qquad
 \rho_b^3=\operatorname{id}_{A_{2,b}}.
 \label{p12-83-c1-s10:eq:u006b-rho-fibra}
\end{equation}
La independencia respecto de \(r\) liga el carácter a la fibra
dodecafásica; \(r\) conserva la posición interna de la vuelta nonádica.

\subsection{Categoría odométrica del TPK: cilindros, truncamientos y morfismos de sucesión}

Recordemos que \(\mathcal O_{9,12}\) tiene objetos
\[
 (b,r)\in\mathbb Z/12\mathbb Z\times\{0,\ldots,8\}.
\]
Para \(n\in\mathbb N\), la flecha única de longitud \(n\) es
\begin{equation}
 \gamma_{(b,r),n}:
 (b,r)\longrightarrow
 \left(
 b+\kappa_r(n),[r+n]_9
 \right),
 \qquad
 \kappa_r(n):=
 \left\lfloor\frac{r+n}{9}\right\rfloor.
 \label{p12-83-c1-s10:eq:u006b-flecha-odometro}
\end{equation}
La composición se apoya en
\begin{equation}
 \kappa_r(n+m)
 =\kappa_r(n)+\kappa_{[r+n]_9}(m).
 \label{p12-83-c1-s10:eq:u006b-cociclo-odometro}
\end{equation}
La longitud \(\ell(\gamma_{(b,r),n})=n\) es aditiva.

\begin{definition}[Funtor ciclotómico de fase--longitud]
\label{p12-83-c1-s10:def:u006b-funtor-ciclotomico}
Sea \(\mathbf{Rep}_{\mathbb C}(C_3)\) la categoría de representaciones
complejas de \(C_3\), y sea \(\mathbf B\mathbb N\) la categoría monoidal
de un solo objeto con endomorfismos \((\mathbb N,+)\). Definimos
\begin{equation}
 \mathfrak F_{\rm cyc}:
 \mathcal O_{9,12}
 \longrightarrow
 \mathbf{Rep}_{\mathbb C}(C_3)\times\mathbf B\mathbb N
 \label{p12-83-c1-s10:eq:u006b-signatura-funtor}
\end{equation}
por
\[
 \mathfrak F_{\rm cyc}(b,r)
 =(A_{2,b}\otimes_{\mathbb Z}\mathbb C,\rho_b)
\]
y
\begin{equation}
 \mathfrak F_{\rm cyc}(\gamma)
 :=
 \left(
 \iota_{b'}^{-1}P_{b'}C^{\ell(\gamma)}
 P_b^{-1}\iota_b,
 \ell(\gamma)
 \right)
 \label{p12-83-c1-s10:eq:u006b-flecha-funtor}
\end{equation}
para \(\gamma:(b,r)\to(b',r')\).
\end{definition}

\begin{theorem}[Functorialidad y separación de los registros]
\label{p12-83-c1-s10:thm:u006b-functorialidad}
La asignación de la \cref{p12-83-c1-s10:def:u006b-funtor-ciclotomico} es un funtor. Su
primera componente conserva únicamente la clase de longitud módulo tres; la
segunda conserva el entero completo y, por reducción, determina esa clase.
La componente ciclotómica no permite reconstruir la longitud absoluta.
\end{theorem}

\begin{proof}
Si \(\gamma_1:(b_0,r_0)\to(b_1,r_1)\) y
\(\gamma_2:(b_1,r_1)\to(b_2,r_2)\) tienen longitudes \(n_1,n_2\), la
trivialización intermedia se cancela:
\begin{align}
 &\bigl(
 \iota_{b_2}^{-1}P_{b_2}C^{n_2}P_{b_1}^{-1}\iota_{b_1}
 \bigr)
 \bigl(
 \iota_{b_1}^{-1}P_{b_1}C^{n_1}P_{b_0}^{-1}\iota_{b_0}
 \bigr)\notag\\
 &\qquad=
 \iota_{b_2}^{-1}P_{b_2}C^{n_1+n_2}
 P_{b_0}^{-1}\iota_{b_0}.
 \label{p12-83-c1-s10:eq:u006b-composicion-funtor}
\end{align}
La segunda componente compone por suma. Como \(C^n\) sólo depende de
\([n]_3\), pero \(\ell=n\) conserva el entero, las dos coordenadas no se
identifican.
\end{proof}

Si
\[
 Q:\mathsf{TPK}\longrightarrow\mathcal O_{9,12}
\]
es el funtor odométrico de U006, la realización tipada es
\begin{equation}
 \mathfrak F_{\rm cyc}\circ Q:
 \mathsf{TPK}\longrightarrow
 \mathbf{Rep}_{\mathbb C}(C_3)\times\mathbf B\mathbb N.
 \label{p12-83-c1-s10:eq:u006b-compuesto-tpk-ciclotomico}
\end{equation}

\subsection{Elevación modular de caminos cardinales}

Sea \(\mathcal A_{N,d}\) la categoría libre de caminos cardinales de
\((\mathbb Z/N\mathbb Z)^d\). Sean
\[
 N\ge2,\quad u\ge1,\quad\gcd(u,N)=1,
 \quad k:=\operatorname{ord}_N(u),
 \quad q:=\frac{u^k-1}{N}.
\]

\begin{definition}[Sustitución cardinal por \(u\)]
\label{p12-83-c1-s10:def:u006b-sustitucion-cardinal}
El funtor
\[
 \mathcal S_u:\mathcal A_{N,d}\longrightarrow\mathcal A_{N,d}
\]
envía un objeto \(p\) a \(up\) y cada arista orientada a la concatenación
ordenada de \(u\) copias con la misma dirección.
\end{definition}

\begin{theorem}[Elevación modular a endomorfismo de caminos]
\label{p12-83-c1-s10:thm:u006b-elevacion-modular}
Para todo objeto \(p\), camino \(\gamma\) y
\[
 \kappa_{N,r}(n):=
 \left\lfloor\frac{r+n}{N}\right\rfloor,
\]
se cumple
\begin{align}
 \mathcal S_u^k(p)&=p,
 \label{p12-83-c1-s10:eq:u006b-retorno-objetos}\\
 \ell(\mathcal S_u^k\gamma)&=u^k\ell(\gamma),
 \label{p12-83-c1-s10:eq:u006b-longitud-sustituida}\\
 \kappa_{N,r}(u^kn)-\kappa_{N,r}(n)&=qn.
 \label{p12-83-c1-s10:eq:u006b-diferencia-vueltas}
\end{align}
\end{theorem}

\begin{proof}
La categoría es libre, de modo que las imágenes de objetos y generadores
determinan un funtor único. Cada aplicación multiplica la longitud por
\(u\). Como \(u^k=1+qN\), los objetos retornan módulo \(N\), y
\[
 \left\lfloor
 \frac{r+n+qNn}{N}
 \right\rfloor
 =
 \left\lfloor\frac{r+n}{N}\right\rfloor+qn.
\]
\end{proof}

En el caso nonádico,
\begin{equation}
 (N,u,k,q)=(9,4,3,7),
 \qquad
 4^3=64=1+7\cdot9.
 \label{p12-83-c1-s10:eq:u006b-parametros-nonadicos}
\end{equation}
Por tanto,
\begin{equation}
 \boxed{
 \mathcal S_4^3(e)
 \text{ conserva los extremos y contiene siete vueltas nonádicas
 orientadas por arista inicial}.}
 \label{p12-83-c1-s10:eq:u006b-siete-vueltas}
\end{equation}
Los extremos y la posición residual retornan en tres aplicaciones; el
estado enriquecido no tiene por qué retornar. Sin una prueba de
invertibilidad, \(\mathcal S_4^3\) es un endomorfismo vertical de caminos,
no una monodromía.

\begin{remark}[Separación respecto de la dilatación del reloj]
\label{p12-83-c1-s10:rem:u006b-s4-no-d4}
\(\mathcal S_4\) sustituye objetos y aristas en una categoría de caminos.
La aplicación \(D_4\) de U006C multiplicará por cuatro la coordenada
\(\Theta_{108}\). Comparten un factor entero, pero no dominio, codominio ni
ley de composición.
\end{remark}

\begin{criterion}[Criterios de refutación]
\label{p12-83-c1-s10:crit:u006b-refutacion}
El bloque queda refutado si falla cualquiera de las identidades
\eqref{p12-83-c1-s10:eq:u006b-identidades-coxeter},
\eqref{p12-83-c1-s10:eq:u006b-cociclo-odometro},
\eqref{p12-83-c1-s10:eq:u006b-composicion-funtor} o
\eqref{p12-83-c1-s10:eq:u006b-parametros-nonadicos}; también queda mal tipado si se
identifican la fase \(C^n\), la longitud \(n\), la sustitución
\(\mathcal S_4\) y la dilatación \(D_4\).
\end{criterion}
}%


% Unidad U004. Redacción autónoma desde los generadors TRIT de 1063/live.

\section{TRIT como unidad mínima de información topológica: orientación ternaria, memoria de acarreo y orden temporal}
\label{p12-83-c1-s11:chap:trit-levantado-regimenes}

La curvatura mixta de APP produjo tres valores orientados: cero en las hojas
puras y signos opuestos al intercambiar el orden suma--producto. El TRIT
tipa esa tríada y conserva, además del residuo visible, el cociente necesario
para reconstruir el entero levantado. Esta unidad separa expresamente cuatro
niveles: clase ternaria, representante balanceado, acarreo y régimen
cuadrático.

\begin{definition}[Unidad de información topológica TRIT]
El \emph{TRIT} es la unidad elemental orientada cuyo espacio de estados es
\(\mathsf{TRIT}=\{-1,0,+1\}\). Su capacidad logarítmica nativa es
\(\log 3\) y su expresión en unidades binarias es \(\log_2 3\). A diferencia
de un bit, sus tres estados no se reducen a una codificación cardinal: tipan,
respectivamente, apertura, umbral y retorno, y conservan el cociente de
acarreo requerido para reconstruir la transición que produjo el residuo
visible.
\end{definition}

\subsection{Sistema ternario orientado y representación balanceada de los estados TRIT}

Definimos el conjunto orientado
\[
 \mathsf{TRIT}:=\{-1,0,+1\}.
\]
La carta balanceada es la biyección de conjuntos
\[
 \operatorname{bal}:\mathbb F_3\longrightarrow\mathsf{TRIT},
 \qquad
 [0]_3\mapsto0,\qquad[1]_3\mapsto+1,\qquad[2]_3\mapsto-1.
\]
No se trata de una igualdad de tipos. \(\mathbb F_3\) porta suma y producto
de cuerpo; \(\mathsf{TRIT}\) porta orden y orientación. Toda operación
transportada entre ambos debe indicar la carta utilizada.

La orientación de las nueve fases es
\[
 \varepsilon_9(r)=
 \begin{cases}
 +1,&r\in\{1,4,7\},\\
 -1,&r\in\{2,5,8\},\\
 0,&r\in\{3,6,9\}.
 \end{cases}
\]
Su tabla completa es
\[
\begin{array}{c|ccccccccc}
r&1&2&3&4&5&6&7&8&9\\\hline
\varepsilon_9(r)&+1&-1&0&+1&-1&0&+1&-1&0
\end{array}.
\]
La periodicidad de tres no elimina la fase nonádica: las fases
\(1,4,7\), por ejemplo, tienen la misma orientación y posiciones distintas
en la vuelta.

\subsection{Algoritmo de división ternaria balanceada y unicidad del par cociente–residuo}

\begin{theorem}[Levantamiento ternario local]
\label{p12-83-c1-s11:thm:trit-levantamiento-local}
Para cada \(n\in\{-4,-3,\ldots,4\}\) existe un único par
\((r,c)\in\mathsf{TRIT}^2\) tal que
\[
 n=r+3c.
\]
\end{theorem}

\begin{proof}
La existencia se verifica en la tabla
\[
\begin{array}{c|rrrrrrrrr}
n&-4&-3&-2&-1&0&1&2&3&4\\\hline
r&-1&0&+1&-1&0&+1&-1&0&+1\\
c&-1&-1&-1&0&0&0&+1&+1&+1
\end{array}.
\]
Si \(r+3c=r'+3c'\), entonces \(r-r'=3(c'-c)\). Como
\(r-r'\in\{-2,-1,0,1,2\}\), el único múltiplo de tres posible es cero;
por tanto \(r=r'\) y \(c=c'\).
\end{proof}

La proyección visible \((r,c)\mapsto r\) posee, sobre el valor cero, tres
estados locales:
\[
 0^-:=(0,-1),
 \qquad
 0^0:=(0,0),
 \qquad
 0^+:=(0,+1).
\]
Publican el mismo residuo, pero reconstruyen \(-3,0,+3\). Esta fibra triple
es la forma mínima de una distinción que reaparecerá en el estado enriquecido
del Topological Phase Kernel.

\begin{corollary}[Iteración balanceada]
\label{p12-83-c1-s11:cor:trit-iteracion-balanceada}
Todo entero \(N\) admite una expansión ternaria balanceada finita
\[
 N=r_0+3r_1+\cdots+3^mr_m,
 \qquad r_j\in\mathsf{TRIT}.
\]
Se obtiene iterando la división \(N_k=r_k+3N_{k+1}\), donde \(r_k\) es el
representante balanceado de \(N_k\bmod3\).
\end{corollary}

\begin{proof}
La división euclídea módulo tres da un residuo en \(\{0,1,2\}\); sustituir
\(2\) por \(-1\) incrementa el cociente en una unidad y produce
\(r_k\in\mathsf{TRIT}\). El valor absoluto del cociente disminuye fuera de
un conjunto finito, por lo que el proceso termina. La unicidad se obtiene
comparando la primera cifra no coincidente y reduciendo módulo tres.
\end{proof}

\subsection{Composición y cociclo de acarreo}

Para \(r,s\in\mathsf{TRIT}\), definimos \(r\oplus s\in\mathsf{TRIT}\) y
\(\kappa_3(r,s)\in\mathbb Z\) por la igualdad única
\[
 r+s=(r\oplus s)+3\kappa_3(r,s).
\]
La tabla es
\[
\begin{array}{c|ccc}
\oplus&-1&0&+1\\\hline
-1&+1&-1&0\\
0&-1&0&+1\\
+1&0&+1&-1
\end{array},
\qquad
\begin{array}{c|ccc}
\kappa_3&-1&0&+1\\\hline
-1&-1&0&0\\
0&0&0&0\\
+1&0&0&+1
\end{array}.
\]

\begin{proposition}[Identidad cocíclica ternaria]
\label{p12-83-c1-s11:prop:trit-cociclo}
Para \(r,s,t\in\mathsf{TRIT}\),
\[
 \kappa_3(r,s)+\kappa_3(r\oplus s,t)
 =\kappa_3(s,t)+\kappa_3(r,s\oplus t).
\]
\end{proposition}

\begin{proof}
Ambos miembros son el cociente total que aparece al escribir
\(r+s+t\) como un representante balanceado más un múltiplo de tres. La
unicidad del levantamiento prueba la igualdad.
\end{proof}

Si \((r,c)\) y \((s,d)\) son estados levantados, su suma es
\[
 (r,c)\boxplus(s,d)
 :=\bigl(r\oplus s,\ c+d+\kappa_3(r,s)\bigr).
\]
La \cref{p12-83-c1-s11:prop:trit-cociclo} prueba que \(\boxplus\) es asociativa y que la
reconstrucción \(\operatorname{rec}_3(r,c)=r+3c\) satisface
\[
 \operatorname{rec}_3(x\boxplus y)
 =\operatorname{rec}_3(x)+\operatorname{rec}_3(y).
\]

\subsection{Involución de orientación y orden temporal}

Para una palabra trítica \(\tau=(\tau_1,\ldots,\tau_n)\), definimos
\[
 \mathcal C_{\rm rev}(\tau_1,\ldots,\tau_n)
 :=(-\tau_n,\ldots,-\tau_1).
\]

\begin{lemma}[Reversión orientada]
\label{p12-83-c1-s11:lem:trit-reversion}
La aplicación \(\mathcal C_{\rm rev}\) es una involución. Intercambia los
sectores \(+1\) y \(-1\) y fija el sector \(0\).
\end{lemma}

\begin{proof}
Aplicarla dos veces invierte dos veces el orden y cambia dos veces cada
signo. Así se recupera la palabra inicial. Las demás afirmaciones son
inmediatas en cada componente.
\end{proof}

Cuando se fija un parámetro temporal orientado, la convención del TPK es
\[
 +1:\text{retorno con memoria},
 \qquad
 0:\text{umbral presente},
 \qquad
 -1:\text{apertura de prolongación}.
\]
Esta asignación no añade un cuarto estado temporal. La dirección se conserva
en el signo y la posición exacta se conserva en la fase nonádica.

\subsection{Familia universal de álgebras cuadráticas}

Para \(\tau\in\mathsf{TRIT}\), definimos
\[
 \mathbb A_\tau:=\mathbb R[J_\tau]/(J_\tau^2+\tau I).
\]
Los tres casos son:
\[
\begin{array}{c|c|c|c}
\tau&J_\tau^2&\text{régimen}&\exp(tJ_\tau)\\\hline
+1&-I&\text{elíptico}&\cos t\,I+\sin t\,J_{+1}\\
0&0&\text{parabólico}&I+tJ_0\\
-1&+I&\text{hiperbólico}&\cosh t\,I+\sinh t\,J_{-1}
\end{array}.
\]

\begin{theorem}[Exponencial cuadrática uniforme]
\label{p12-83-c1-s11:thm:trit-exponencial-uniforme}
Si \(J_\tau^2=-\tau I\), entonces
\[
 \exp(tJ_\tau)
 =C_\tau(t)I+S_\tau(t)J_\tau,
\]
donde
\[
\begin{array}{c|ccc}
\tau&+1&0&-1\\\hline
C_\tau(t)&\cos t&1&\cosh t\\
S_\tau(t)&\sin t&t&\sinh t
\end{array}.
\]
\end{theorem}

\begin{proof}
Separamos la serie exponencial en potencias pares e impares. Las relaciones
\(J_{+1}^{2m}=(-1)^mI\), \(J_0^2=0\) y
\(J_{-1}^{2m}=I\) producen, respectivamente, las series de coseno y seno,
la truncación lineal y las series hiperbólicas.
\end{proof}

Sobre \(\mathbb F_3\) aparecen los análogos
\[
 \mathbb F_9=\mathbb F_3[\iota]/(\iota^2+1),
 \qquad
 \mathbb D_3=\mathbb F_3[\epsilon]/(\epsilon^2),
 \qquad
 \mathbb S_3=\mathbb F_3[j]/(j^2-1).
\]
El primero es un cuerpo de nueve elementos; el segundo contiene un nilpotente
no nulo; el tercero se descompone mediante los idempotentes
\((1\pm j)/2\) y es isomorfo a \(\mathbb F_3\times\mathbb F_3\). La
coincidencia cardinal no identifica estas tres álgebras.

La fuente integral del régimen elíptico es
\[
 \mathbb Q_{\mathbb Z}:=\mathbb Z[u]/(u^2+1).
\]
La reducción módulo tres envía \(u\) a \(\iota\) y produce
\(\mathbb F_9\). La extensión de escalares a \(\mathbb R\) y una
representación \(u\mapsto J_{+1}\) producen el realizador real elíptico.
No existe por ello un mapa no tipado \(\mathbb F_9\to\mathbb R\): ambas
realizaciones proceden de una fuente integral común.

\subsection{Representaciones matriciales del TRIT seleccionadas por el signo de \(J_\tau^2\)}

Sea
\[
 C=\begin{pmatrix}0&-1\\1&-1\end{pmatrix},
 \qquad C^3=I,
\]
y definamos
\[
 J_e:=\frac{2C+I}{\sqrt3}
 =\frac1{\sqrt3}\begin{pmatrix}1&-2\\2&-1\end{pmatrix}.
\]
Una multiplicación directa da
\[
 (2C+I)^2=-3I,
 \qquad J_e^2=-I.
\]
El realizador parabólico es
\[
 J_p:=N=\begin{pmatrix}0&1\\0&0\end{pmatrix},
 \qquad N^2=0.
\]
Para el régimen hiperbólico tomamos
\[
 F_{\rm av}=\begin{pmatrix}0&1\\1&1\end{pmatrix},
 \qquad
 A:=F_{\rm av}^2=\begin{pmatrix}1&1\\1&2\end{pmatrix},
\]
y
\[
 J_h:=\frac{2A-3I}{\sqrt5}
 =\frac1{\sqrt5}\begin{pmatrix}-1&2\\2&1\end{pmatrix}.
\]
Entonces
\[
 (2A-3I)^2=5I,
 \qquad J_h^2=I.
\]

La forma bilineal
\[
 G=\begin{pmatrix}2&1\\1&-2\end{pmatrix}
\]
satisface
\[
 A^TGA=G.
\]
Además,
\[
 \cosh(2\log\varphi)=\frac32,
 \qquad
 \sinh(2\log\varphi)=\frac{\sqrt5}{2},
\]
por lo que
\[
 A=\exp(2\log\varphi\,J_h)
  =\cosh(2\log\varphi)I+\sinh(2\log\varphi)J_h.
\]
Esta identidad será recuperada causalmente en la biografía de
\(\varphi\); aquí sólo se verifica la representación matricial.

Los tres controles exponenciales son
\[
 \exp(\pi J_e)+I=0,
 \qquad
 \exp(J_p)=I+J_p,
 \qquad
 \exp(2\log\varphi\,J_h)=F_{\rm av}^2.
\]
La presencia de \(\pi\) y \(\varphi\) en estas identidades de verificación
no autoriza a usarlas como constructores de dichas constantes. Sus
constructores aparecerán después desde la genealogía nonádica.

\subsection{Acoplamiento local APP--TRIT}

Con el intercambio
\[
 R=\begin{pmatrix}0&1\\1&0\end{pmatrix}
\]
definimos, para \(\tau\in\mathsf{TRIT}\),
\[
 B_\tau=(4+3\tau)I+(5-3\tau)R.
\]
Puesto que \(P_\pm=(I\pm R)/2\),
\[
 B_\tau=9P_++(6\tau-1)P_-.
\]

\begin{theorem}[Recuperación espectral de la orientación]
\label{p12-83-c1-s11:thm:app-trit-recuperacion-espectral}
El espectro de \(B_\tau\) es
\[
 \operatorname{Spec}(B_\tau)=\{9,6\tau-1\}.
\]
El autovalor diferencial \(\lambda_-\) determina unívocamente el TRIT:
\[
 \tau=\frac{\lambda_-+1}{6}.
\]
Además, el bloque central de APP es \(B_{+1}\).
\end{theorem}

\begin{proof}
Los proyectores \(P_+\) y \(P_-\) diagonalizan \(R\) con autovalores
\(+1\) y \(-1\). La fórmula espectral se obtiene sustituyendo. La segunda
identidad despeja \(\tau\). Para \(\tau=+1\),
\(B_{+1}=7I+2R=\begin{psmallmatrix}7&2\\2&7\end{psmallmatrix}\), que es
el bloque único de la \cref{p12-83-c1-s1:thm:app-bloque-central}.
\end{proof}

Por tanto, el TRIT no se añade como etiqueta retrospectiva: puede recuperarse
del modo diferencial de una familia construida sobre el intercambio interno
de APP.

\subsection{Obstrucciones del TRIT: separación residuo–representante–acarreo, involutividad y recuperación espectral de \(\tau\)}

\begin{criterion}[Criterios de refutación del TRIT]
La unidad queda refutada si:
\begin{enumerate}
 \item se identifican clase residual, representante balanceado y acarreo;
 \item se afirma que \(0^-\), \(0^0\) y \(0^+\) son el mismo estado;
 \item la composición \(\boxplus\) deja de reconstruir la suma entera;
 \item \(\mathcal C_{\rm rev}\) no es involutiva;
 \item se identifican las tres álgebras por tener el mismo cardinal;
 \item algún realizador incumple \(J_\tau^2=-\tau I\);
 \item el valor de \(\tau\) se elige después de observar el resultado y no
       se recupera del espectro de \(B_\tau\).
\end{enumerate}
\end{criterion}

APP proporciona el soporte, la matriz multiplicativa y la ley aditiva
interna; TRIT proporciona orientación,
levantamiento y tipo cuadrático. La unidad siguiente construirá el
Topological Phase Kernel como categoría de transporte de esos datos, sin
confundir el selector interno con el estado completo.

\HMTDeferredTPKBase
\HMTDeferredTPKDevelopment


% Unidad U006C. Dilatación aritmética, cociclo de acarreo y
% semiconjugación cuaternaria.

\section{Dilatación cuaternaria del registro, obstrucción observable y semiconjugación refinada}
\label{p12-83-c1-s12:chap:dilatacion-cuaternaria-semiconjugacion}

La sustitución cardinal \(\mathcal S_4\) construida en la unidad anterior
actúa sobre caminos. En esta unidad se define una segunda operación, de
naturaleza aritmética, sobre el registro
\(\mathbb Z/12\mathbb Z\times\mathbb Z/9\mathbb Z\), y una tercera
operación, categorial, que subdivide cada transición observable en cuatro
microtransiciones. El objetivo no es fusionarlas, sino demostrar la relación
exacta entre sus realizaciones.

\subsection{Dilatación sobre el registro producto}

Sea
\[
 \mathcal Q_{12,9}:=
 \mathbb Z/12\mathbb Z\times\mathbb Z/9\mathbb Z.
\]
Para \(r\in\mathbb Z/9\mathbb Z\), denotamos por
\(\widetilde r\in\{0,\ldots,8\}\) su representante normalizado.

\begin{definition}[Dilatación cuaternaria del registro]
\label{p12-83-c1-s12:def:u006c-dilatacion-registro}
Definimos
\begin{equation}
 D_4(b,r):=
 \left(
 \left[
  4b+\left\lfloor\frac{4\widetilde r}{9}\right\rfloor
 \right]_{12},
 [4r]_9
 \right).
 \label{p12-83-c1-s12:eq:u006c-d4-producto}
\end{equation}
La coordenada total es
\begin{equation}
 \Theta:\mathcal Q_{12,9}\xrightarrow{\;\sim\;}
 \mathbb Z/108\mathbb Z,
 \qquad
 \Theta(b,r):=[9b+\widetilde r]_{108}.
 \label{p12-83-c1-s12:eq:u006c-coordenada-total}
\end{equation}
\end{definition}

\begin{proposition}[Realización por multiplicación modular]
\label{p12-83-c1-s12:prop:u006c-d4-total}
En la coordenada total,
\begin{equation}
 \boxed{
 \Theta\circ D_4\circ\Theta^{-1}(\theta)=[4\theta]_{108}.}
 \label{p12-83-c1-s12:eq:u006c-d4-total}
\end{equation}
En particular,
\begin{equation}
 \operatorname{im}(D_4)
 =4\mathbb Z/108\mathbb Z
 \simeq\mathbb Z/27\mathbb Z,
 \label{p12-83-c1-s12:eq:u006c-imagen-d4}
\end{equation}
y la restricción inducida sobre la imagen tiene orden nueve.
\end{proposition}

\begin{proof}
Para \(\theta=9b+\widetilde r\),
\[
 4\theta
 =9\left(
  4b+\left\lfloor\frac{4\widetilde r}{9}\right\rfloor
 \right)
 +(4\widetilde r\bmod9).
\]
Ésta es exactamente la descomposición cociente--residuo de
\eqref{p12-83-c1-s12:eq:u006c-d4-producto}. La imagen está formada por las clases
múltiplos de cuatro. Tras dividir por cuatro se obtiene
\(\mathbb Z/27\mathbb Z\), y
\[
\begin{aligned}
 4^1&\equiv4,& 4^2&\equiv16,& 4^3&\equiv10,\\
 4^4&\equiv13,& 4^5&\equiv25,& 4^6&\equiv19,\\
 4^7&\equiv22,& 4^8&\equiv7,& 4^9&\equiv1\pmod{27}.
\end{aligned}
\]
Ninguna potencia anterior es uno; por tanto
\(\operatorname{ord}_{27}(4)=9\).
\end{proof}

\subsection{Identidad exacta de acarreo}

Para \(r\in\{0,\ldots,8\}\) y \(n\in\mathbb N\), definimos
\[
 \kappa_r(n):=\left\lfloor\frac{r+n}{9}\right\rfloor,
 \qquad
 r':=[r+n]_9.
\]

\begin{theorem}[Cociclo cuaternario de acarreo]
\label{p12-83-c1-s12:thm:u006c-cociclo-d4}
Se cumple
\begin{equation}
 \boxed{
 \kappa_{[4r]_9}(4n)
 =4\kappa_r(n)
 +\left\lfloor\frac{4\widetilde r'}{9}\right\rfloor
 -\left\lfloor\frac{4\widetilde r}{9}\right\rfloor.}
 \label{p12-83-c1-s12:eq:u006c-cociclo-d4}
\end{equation}
La última diferencia es el coborde producido por la elección de
representantes \(r\mapsto\widetilde r\).
\end{theorem}

\begin{proof}
Escribamos
\[
 r+n=9\kappa_r(n)+\widetilde r'.
\]
Al multiplicar por cuatro,
\[
 4r+4n=36\kappa_r(n)+4\widetilde r'.
\]
Separar cociente y residuo módulo nueve y restar el cociente ya contenido en
\(4\widetilde r\) produce
\eqref{p12-83-c1-s12:eq:u006c-cociclo-d4}.
\end{proof}

\begin{corollary}[Tercera iteración]
\label{p12-83-c1-s12:cor:u006c-tercera-iteracion}
Para todo \((b,r)\in\mathcal Q_{12,9}\),
\begin{equation}
 \boxed{D_4^3(b,r)=([4b+7\widetilde r]_{12},r).}
 \label{p12-83-c1-s12:eq:u006c-d4-cubo}
\end{equation}
\end{corollary}

\begin{proof}
Por la \cref{p12-83-c1-s12:prop:u006c-d4-total},
\[
 \Theta(D_4^3(b,r))
 =[4^3(9b+\widetilde r)]_{108}
 =[64(9b+\widetilde r)]_{108}.
\]
Como \(64=1+7\cdot9\), el residuo módulo nueve vuelve a ser \(r\), mientras
la coordenada de bloque se transforma en
\([4b+7\widetilde r]_{12}\).
\end{proof}

El retorno de \(r\) no es retorno del registro completo. La diferencia
\(7\widetilde r\) conserva la memoria del representante interno atravesado.

\subsection{Obstrucción en el cociente observable}

En la configuración observable de
\cref{p12-83-c1-s3:def:tpk-configuracion-observable}, escribimos
\[
 q=(p^\Sigma,p^\Pi,d^\Sigma,d^\Pi,\theta)
 \in\mathcal Q_{\rm obs}.
\]
La proyección posicional es
\begin{equation}
 \pi_{\rm pos}(q):=(p^\Sigma,p^\Pi)\in X_9^2.
 \label{p12-83-c1-s12:eq:u006c-proyeccion-posicional}
\end{equation}
La firma temporal es
\[
 \tau(q):=\varepsilon_9(r(\theta))\in\{-1,0,+1\}.
\]
Definimos el desplazamiento observable orientado
\begin{equation}
 a(q):=
 \begin{cases}
  (\delta(d^\Sigma),0),&\tau(q)=+1,\\
  (0,\delta(d^\Pi)),&\tau(q)=-1,\\
  (0,0),&\tau(q)=0.
 \end{cases}
 \label{p12-83-c1-s12:eq:u006c-desplazamiento-observable}
\end{equation}
Entonces
\[
 \pi_{\rm pos}(F_{\rm obs}q)=\pi_{\rm pos}(q)+a(q),
\]
con la actualización de direcciones y reloj definida en U005.

\begin{theorem}[Obstrucción al endofuntor observable ordinario]
\label{p12-83-c1-s12:thm:u006c-obstruccion}
No existe un endofuntor del espacio observable ordinario que satisfaga
simultáneamente:
\begin{enumerate}
 \item multiplicar la posición geométrica por cuatro;
 \item sustituir cada transición por cuatro microtransiciones genuinas;
 \item conservar hoja y orientación en cada microtransición;
 \item enviar una fibra dodecafásica completa a una sola fibra
 dodecafásica.
\end{enumerate}
La obstrucción persiste en un paso de umbral con \(a(q)=0\).
\end{theorem}

\begin{proof}
Sea
\[
 \mathcal B_b:=\{(b,r):0\leq r\leq8\}.
\]
Sustituir una iteración elemental por cuatro fuerza
\[
 h(\theta+1)=h(\theta)+4
\]
y, por inducción,
\[
 h(\theta)=4\theta+c\pmod{108}.
\]
Tomemos \(0\leq\widetilde c<108\) y escribamos
\(\widetilde c=9c_b+c_0\), con \(0\leq c_0<9\). Si toda la fibra
\(\mathcal B_b\) fuese enviada a una sola fibra, el índice de bloque de
destino tendría que ser independiente de \(r\); sin embargo contiene
\begin{equation}
 4b+c_b+
 \left\lfloor\frac{4r+c_0}{9}\right\rfloor
 \pmod{12},
 \label{p12-83-c1-s12:eq:u006c-bloque-obstruccion}
\end{equation}
que no es constante en \(r\). Además, el operador de transición neutra
satisface \(P_0^{\rm TPK}=I\), distinto de las trivializaciones
ciclotómicas \(P_b\); el calendario ordinario \((+1,-1,0)\) hace que
cuatro iteraciones elementales visiten ambos modos aritméticos y el modo
neutro.
\end{proof}

\subsection{Refinamiento mínimo por microinstantes}

\begin{definition}[Estado observable subdividido]
\label{p12-83-c1-s12:def:u006c-estado-subdividido}
Sea
\begin{equation}
 \widetilde{\mathcal Q}_4
 :=\mathcal Q_{\rm obs}\times\{0,1,2,3\},
 \qquad
 \widetilde\pi(q,j)
 :=4\pi_{\rm pos}(q)+j\,a(q).
 \label{p12-83-c1-s12:eq:u006c-estado-refinado}
\end{equation}
Definimos
\begin{equation}
 \widetilde F(q,j):=
 \begin{cases}
  (q,j+1),&0\leq j<3,\\
  (F_{\rm obs}q,0),&j=3,
 \end{cases}
 \qquad
 \iota_4(q):=(q,0).
 \label{p12-83-c1-s12:eq:u006c-dinamica-refinada}
\end{equation}
\end{definition}

\begin{theorem}[Semiconjugación a cuatro pasos]
\label{p12-83-c1-s12:thm:u006c-semiconjugacion}
La inclusión \(\iota_4\) satisface
\begin{equation}
 \boxed{
 \widetilde F^4\circ\iota_4
 =\iota_4\circ F_{\rm obs}.}
 \label{p12-83-c1-s12:eq:u006c-semiconjugacion}
\end{equation}
Las posiciones intermedias son exactamente
\[
 4\pi_{\rm pos}(q)+j\,a(q),
 \qquad j=0,1,2,3.
\]
\end{theorem}

\begin{proof}
Desde \((q,0)\), las primeras tres aplicaciones incrementan sólo el índice
de microinstante. La cuarta aplica \(F_{\rm obs}\) y devuelve el índice a
cero. Esto da
\[
 (q,0)\mapsto(q,1)\mapsto(q,2)\mapsto(q,3)
 \mapsto(F_{\rm obs}q,0).
\]
La fórmula de las posiciones es la definición de
\(\widetilde\pi\). El refinamiento hace visible una subdivisión determinada;
no añade una decisión dinámica.
\end{proof}

\subsection{Categoría subdividida y elevación ciclotómica}

Sea \(\operatorname{sd}_4\mathsf P_{\rm obs}\) la categoría obtenida al
reemplazar cada generador
\(\gamma:q\to F_{\rm obs}q\) por la cadena
\begin{equation}
 (q,0)\to(q,1)\to(q,2)\to(q,3)\to(F_{\rm obs}q,0).
 \label{p12-83-c1-s12:eq:u006c-cadena-subdivision}
\end{equation}
El funtor
\[
 \operatorname{sd}_4:
 \mathsf P_{\rm obs}\longrightarrow
 \operatorname{sd}_4\mathsf P_{\rm obs}
\]
envía cada generador a esa cadena. El funtor de colapso
\[
 \operatorname{col}_4:
 \operatorname{sd}_4\mathsf P_{\rm obs}
 \longrightarrow\mathsf P_{\rm obs}
\]
recompone las cuatro microflechas y satisface
\begin{equation}
 \operatorname{col}_4\circ\operatorname{sd}_4
 =\operatorname{id}_{\mathsf P_{\rm obs}}.
 \label{p12-83-c1-s12:eq:u006c-colapso}
\end{equation}

En una cadena que permanece en la fibra \(b\), las tres primeras
microflechas actúan por \(\rho_b\); la cuarta incorpora el transporte
interfibra si existe cruce dodecafásico. Como \(\rho_b^3=I\), el producto
ordenado de las cuatro acciones coincide con el morfismo ciclotómico de la
flecha original.

\begin{proposition}[Invariancia ciclotómica por subdivisión]
\label{p12-83-c1-s12:prop:u006c-invariancia-ciclotomica}
Existe un levantamiento tipado
\[
 \widetilde{\mathfrak F}_{\rm cyc}:
 \operatorname{sd}_4\mathsf P_{\rm obs}
 \longrightarrow
 \mathbf{Rep}_{\mathbb C}(C_3)\times\mathbf B\mathbb N
\]
tal que, para toda flecha \(\gamma\) de
\(\mathsf P_{\rm obs}\),
\begin{equation}
 \boxed{
 \widetilde{\mathfrak F}_{\rm cyc}
 \bigl(\operatorname{sd}_4\gamma\bigr)
 =
 \bigl(\mathfrak F_{\rm cyc}\circ Q\bigr)(\gamma).}
 \label{p12-83-c1-s12:eq:u006c-invariancia-ciclotomica}
\end{equation}
\end{proposition}

\begin{proof}
Se asignan a las microflechas internas las potencias sucesivas de
\(\rho_b\), y a la cuarta el transporte entre los marcos de las fibras de
origen y destino. La composición de los tres primeros factores es la
identidad porque \(\rho_b^3=I\). El factor restante es exactamente el
morfismo que \(\mathfrak F_{\rm cyc}\circ Q\) asigna a la flecha
colapsada. La categoría subdividida es libre sobre esas microflechas, por lo
que la asignación se extiende de manera única.
\end{proof}

\subsection{\(3\) operaciones cuaternarias no intercambiables}

La arquitectura contiene tres operaciones distintas:
\begin{enumerate}
 \item \(\mathcal S_4\), sustitución cardinal de objetos y caminos;
 \item \(D_4\), multiplicación por cuatro en
 \(\Theta_{108}\);
 \item \(\operatorname{sd}_4\), subdivisión functorial de una transición
 observable.
\end{enumerate}
Sus compatibilidades están dadas, respectivamente, por
\[
 4^3=1+7\cdot9,\qquad
 \Theta D_4=4\Theta,\qquad
 \widetilde F^4\iota_4=\iota_4F_{\rm obs}.
\]
Compartir el entero cuatro no convierte ninguna en otra.

\begin{criterion}[Criterios de refutación]
\label{p12-83-c1-s12:crit:u006c-refutacion}
La construcción queda refutada si falla cualquiera de las identidades
\eqref{p12-83-c1-s12:eq:u006c-d4-total},
\eqref{p12-83-c1-s12:eq:u006c-cociclo-d4},
\eqref{p12-83-c1-s12:eq:u006c-d4-cubo},
\eqref{p12-83-c1-s12:eq:u006c-semiconjugacion} o
\eqref{p12-83-c1-s12:eq:u006c-invariancia-ciclotomica}; también si una realización
observable satisface las cuatro condiciones del
\cref{p12-83-c1-s12:thm:u006c-obstruccion} sin conservar un registro equivalente al
microinstante.
\end{criterion}


% Unidad U006. Lectores, odómetro y calendario; redacción autónoma.

\section{Lectores modulares, odómetro y calendario dodecafásico}
\label{p12-83-c1-s13:chap:tpk-lectores-odometro-calendario}

El estado enriquecido del TPK no se identifica con ninguna de sus lecturas.
Esta unidad define los lectores módulo nueve, módulo tres y módulo mil, el
odómetro fase--sector, el cociente predictivo de treinta y seis estados y la
extensión cíclica de ciento ocho posiciones. Cada construcción declara qué
publica y qué memoria conserva.

\subsection{Esquema tipado de los lectores modulares del TPK: dominio, codominio, publicación y memoria conservada}

\begin{definition}[Ficha operatoria]
\label{p12-83-c1-s13:def:tpk-ficha-operatoria}
Un operador canónico del TPK se registra mediante una séxtupla
\[
 \mathcal O=(D_{\mathcal O},C_{\mathcal O},f_{\mathcal O},
 I_{\mathcal O},L_{\mathcal O},M_{\mathcal O}),
\]
donde \(D_{\mathcal O}\) es el dominio, \(C_{\mathcal O}\) el codominio,
\(f_{\mathcal O}\) la regla, \(I_{\mathcal O}\) los invariantes preservados,
\(L_{\mathcal O}\) la información publicada u olvidada y
\(M_{\mathcal O}\) la memoria que permanece en el estado enriquecido.
\end{definition}

Dos símbolos con igual cardinal o igual periodo no representan el mismo
operador salvo que exista una equivalencia que entrelace sus dominios,
reglas e invariantes. Esta convención impedirá confundir el cociente
predictivo \(C_{36}^{\rm pred}\), la frontera \(R_{36}\) y un defecto
cardinal de valor treinta y seis.

\subsection{Lector nonádico del estado enriquecido: fase, transición y retorno con memoria}

Para \(n\geq1\), definimos
\[
 \rho_9(n)=1+((n-1)\bmod9),
 \qquad
 q_9(n)=\left\lfloor\frac{n-1}{9}\right\rfloor.
\]
Entonces
\[
 n=9q_9(n)+\rho_9(n).
\]
La ficha de \(\rho_9\) es
\[
 \bigl(
 \mathbb Z_{>0},\mathcal D_9,\rho_9,
 n\equiv\rho_9(n)\pmod9,
 \text{fase positiva},
 q_9(n)
 \bigr).
\]
La raíz digital positiva coincide numéricamente con \(\rho_9\), pero su
semántica es distinta: una se presenta como sección del cociente; la otra,
como suma iterada de cifras.

\subsection{Lector trítico del estado orientado: orden temporal y transporte de acarreo}

La división balanceada se escribe
\[
 n=r_3(n)+3c_3(n),
 \qquad r_3(n)\in\mathsf{TRIT},\quad c_3(n)\in\mathbb Z.
\]
Su ficha es
\[
 \bigl(
 \mathbb Z,\mathsf{TRIT}\times\mathbb Z,
 n\mapsto(r_3(n),c_3(n)),
 \operatorname{rec}_3(r,c)=r+3c,
 r_3,
 c_3
 \bigr).
\]
La publicación puede mostrar sólo \(r_3\), pero el TPK conserva
\(c_3\) si una composición posterior requiere reconstrucción.

\subsection{Lector posicional por bloques en base \(10^3\)}

La división euclídea define
\[
 N=1000q_{1000}(N)+r_{1000}(N),
 \qquad 0\leq r_{1000}(N)<1000.
\]
Su ficha es
\[
 \bigl(
 \mathbb Z,\mathbb Z\times\{0,\ldots,999\},
 N\mapsto(q_{1000},r_{1000}),
 \operatorname{rec}_{1000}(q,r)=1000q+r,
 r_{1000},
 q_{1000}
 \bigr).
\]

\begin{proposition}[Concatenación y lectura residual]
\label{p12-83-c1-s13:prop:tpk-concatenacion-base-mil}
Si \(u_1,\ldots,u_m\in\{0,\ldots,999\}\) y
\[
 D_m=1000D_{m-1}+u_m,
 \qquad D_0=0,
\]
entonces
\[
 D_m=\sum_{j=1}^m u_j1000^{m-j}.
\]
Además,
\[
 D_m\equiv\sum_{j=1}^m u_j\pmod9
\]
porque \(1000\equiv1\pmod9\).
\end{proposition}

\begin{proof}
La primera igualdad se obtiene por inducción en \(m\). La segunda procede
al reducirla módulo nueve.
\end{proof}

La congruencia no reconstruye el orden de los bloques. El registro genealógico conserva la
secuencia \((u_1,\ldots,u_m)\); el lector residual sólo publica su suma
módulo nueve.

\subsection{Categoría odométrica del TPK: objetos cilíndricos y composición por truncamiento}

\begin{definition}[Odómetro nonádico--dodecafásico]
\label{p12-83-c1-s13:def:tpk-odometro-9-12}
La categoría \(\mathcal O_{9,12}\) tiene objetos
\[
 (b,r)\in C_{12}\times\{0,\ldots,8\}.
\]
Una flecha de longitud \(n\in\mathbb N\) es
\[
 (b,r)\xrightarrow{\ n\ }
 \left(
  b+\left\lfloor\frac{r+n}{9}\right\rfloor,
  r+n\bmod9
 \right).
\]
La composición suma longitudes.
\end{definition}

Sea
\[
 \kappa_r(n)=\left\lfloor\frac{r+n}{9}\right\rfloor.
\]

\begin{lemma}[Cociclo del odómetro]
\label{p12-83-c1-s13:lem:tpk-cociclo-odometro}
Si \(r_1=r+n_1\bmod9\), entonces
\[
 \kappa_r(n_1+n_2)
 =\kappa_r(n_1)+\kappa_{r_1}(n_2).
\]
\end{lemma}

\begin{proof}
Escribimos \(r+n_1=9\kappa_r(n_1)+r_1\), con
\(0\leq r_1<9\). Al añadir \(n_2\) y dividir de nuevo entre nueve se
obtiene la identidad.
\end{proof}

Para un estado \(x\) cuya base contiene \(\theta\), definimos
\[
 Q_0(x)=(\beta(\theta),\widetilde\theta\bmod9).
\]
Una flecha TPK tiene longitud igual al número de generadores observables que
la componen.

\begin{theorem}[Funtor fase--longitud]
\label{p12-83-c1-s13:thm:tpk-funtor-fase-longitud}
La asignación que envía un estado \(x\) a \(Q_0(x)\) y una flecha de
longitud \(n\) a la flecha odométrica de longitud \(n\) define un funtor
\[
 Q:\mathsf{TPK}\longrightarrow\mathcal O_{9,12}.
\]
\end{theorem}

\begin{proof}
Cada generador aumenta \(\theta\) en una unidad. Tras \(n\) pasos, el
residuo de fase aumenta en \(n\) módulo nueve y el sector aumenta en
\(\lfloor(r+n)/9\rfloor\). La identidad tiene longitud cero y el
\cref{p12-83-c1-s13:lem:tpk-cociclo-odometro} prueba la compatibilidad con la composición.
\end{proof}

Después de nueve pasos,
\[
 (b,r)\longmapsto(b+1,r).
\]
Después de ciento ocho,
\[
 (b,r)\longmapsto(b,r).
\]
Estas identidades pertenecen al cociente odométrico. No afirman que hayan
retornado cilindro, frontera, registro genealógico o memoria vertical.

\subsection{Cociente mínimo de memoria predictiva}

Sea
\[
 X_{9,12}=\{(a,m):a\in\mathbb Z/9\mathbb Z, 0\leq m<12\}
\]
con sucesor
\[
 S(a,m)=
 \begin{cases}
  (a,m+1),&m<11,\\
  (a+1\bmod9,0),&m=11.
 \end{cases}
\]
Definimos los observables
\[
 \beta_{\rm pred}(a,m)=4a\pmod{12},
 \qquad
 d_{\rm pred}(a,m)=1+(m+\beta_{\rm pred}(a,m)\pmod{12}).
\]
Para un observable \(q\), la palabra futura inclusiva de horizonte \(h\)
es
\[
 Q_h^q(x)=(q(x),q(Sx),\ldots,q(S^hx)).
\]

\begin{theorem}[Cociente predictivo mínimo]
\label{p12-83-c1-s13:thm:tpk-cociente-predictivo-36}
Los observables \(\beta_{\rm pred}\), \(d_{\rm pred}\) y el par
\((\beta_{\rm pred},d_{\rm pred})\) producen el mismo cociente estable
\[
 \pi_{36}(a,m)=(a\bmod3,m).
\]
Tiene treinta y seis estados y fibras de cardinal tres. Los horizontes
mínimos son \(11\), \(8\) y \(0\), respectivamente.
\end{theorem}

\begin{proof}
Para un sistema finito \((X,T)\), sea
\[
 E_h=\operatorname{Eq}(Q_h^q).
\]
Entonces
\[
 E_{h+1}=E_h\cap(T\times T)^{-1}(E_h).
\]
La cadena decreciente estabiliza y su límite es la mayor equivalencia
\(T\)-compatible contenida en \(\operatorname{Eq}(q)\). En el sistema
declarado, la enumeración exacta produce el número de clases siguiente:
\[
\begin{array}{c|rrrrrrrrrrrr}
h&0&1&2&3&4&5&6&7&8&9&10&11\\\hline
\#(X/E_h^{\beta})&3&6&9&12&15&18&21&24&27&30&33&36
\end{array},
\]
y
\[
\begin{array}{c|rrrrrrrrr}
h&0&1&2&3&4&5&6&7&8\\\hline
\#(X/E_h^{d})&12&15&18&21&24&27&30&33&36.
\end{array}
\]
El par ya distingue las treinta y seis clases en horizonte cero. En todos
los casos, las clases estables conservan \(m\) y \(a\bmod3\); las tres
posibilidades de \(a\) con igual residuo módulo tres forman cada fibra.
\end{proof}

Denotamos este objeto por \(C_{36}^{\rm pred}\). No es la frontera
\(R_{36}\), que contendrá tres canales ordenados de seis trits, ni un censo
de treinta y seis elementos de otra construcción.

\subsection{Extensión central del calendario}

Definimos
\[
 \iota:C_{12}\longrightarrow C_{108},
 \qquad \iota([b]_{12})=[9b]_{108},
\]
y
\[
 p:C_{108}\longrightarrow C_9,
 \qquad p([\theta]_{108})=[\theta]_9.
\]
Se obtiene la sucesión exacta
\begin{equation}
 0\longrightarrow C_{12}
 \xrightarrow{\ \iota\ }C_{108}
 \xrightarrow{\ p\ }C_9
 \longrightarrow0.
\label{p12-83-c1-s13:eq:tpk-extension-12-108-9}
\end{equation}
Una sección de conjuntos \(s([r]_9)=[r]_{108}\),
\(0\leq r<9\), induce el cociclo
\[
 c(r,r')=
 \left[\left\lfloor\frac{r+r'}{9}\right\rfloor\right]_{12}.
\]

\begin{proposition}[La extensión no se escinde]
\label{p12-83-c1-s13:prop:tpk-extension-no-escindida}
La sucesión \cref{p12-83-c1-s13:eq:tpk-extension-12-108-9} no es isomorfa a la extensión
producto \(C_{12}\times C_9\).
\end{proposition}

\begin{proof}
El exponente de \(C_{108}\) es \(108\). El exponente de
\(C_{12}\times C_9\) es
\(\operatorname{mcm}(12,9)=36\). Dos grupos finitos isomorfos tienen el
mismo exponente; por tanto, no son isomorfos.
\end{proof}

\begin{theorem}[Necesidad y periodo exacto de ciento ocho posiciones]
\label{p12-83-c1-s13:thm:tpk-periodo-exacto-108}
La traslación \(\theta\mapsto\theta+1\) sobre \(C_{108}\) tiene orden
exactamente \(108\). En consecuencia, una iteración \(n\) retorna el
calendario completo si y sólo si \(108\mid n\).
\end{theorem}

\begin{proof}
La \(n\)-ésima potencia envía \([\theta]\) a \([\theta+n]_{108}\). Es
la identidad para todo \(\theta\) exactamente cuando \([n]_{108}=0\), es
decir, cuando \(108\mid n\).
\end{proof}

En particular, \(1008\) no es otro retorno del calendario:
\[
 1008\equiv36\pmod{108}.
\]
Tras \(1008\) pasos, la fase módulo nueve retorna, pero la coordenada de
\(C_{108}\) se ha desplazado treinta y seis posiciones. Naturalmente,
\(216,324,\ldots\) también retornan el calendario, pero son repeticiones de
la vuelta fundamental de longitud \(108\).

\subsection{Memoria vertical no acotada}

Sea \(q_{\rm mem}\in\mathbb Z_{\geq0}\) el número no acotado de vueltas
nonádicas y sea
\[
 \beta=[q_{\rm mem}]_{12}.
\]
Entonces una vuelta de nueve pasos satisface
\[
 q_{\rm mem}\longmapsto q_{\rm mem}+1,
 \qquad
 \beta\longmapsto\beta+1\pmod{12}.
\]
Después de ciento ocho pasos,
\[
 \beta\longmapsto\beta,
 \qquad
 q_{\rm mem}\longmapsto q_{\rm mem}+12.
\]
El calendario retorna; la memoria vertical no.

\begin{corollary}[Jerarquía de retornos]
\label{p12-83-c1-s13:cor:tpk-jerarquia-retornos}
En el cociente de fase,
\(9\) pasos constituyen una vuelta. Los morfismos compuestos
\[
 \mathcal K_{27}=F_{\rm obs}^{27},
 \qquad
 F_{\rm obs}^{54},
 \qquad
 F_{\rm obs}^{108}
\]
contienen, respectivamente, tres, seis y doce vueltas nonádicas. Ninguno
borra por definición el registro genealógico o la frontera.
\end{corollary}

\subsection{Sistema mínimo de tipos del estado enriquecido y compatibilidades entre sus lectores}

La tabla siguiente responde de manera explícita a qué objetos internos
existen ya en el TPK antes de introducir semillas:
\[
\begin{array}{p{.23\textwidth}p{.27\textwidth}p{.38\textwidth}}
\toprule
Objeto u operador&Dominio y codominio&Función\\\midrule
\(X_9\)&\((\mathbb Z/9)^2\)&soporte APP\\
\(T_N,T_E,T_S,T_O\)&\(X_9\to X_9\)&traslaciones cardinales\\
\(M_{\rm ph}\)&\(\mathcal D_9\to\mathsf{TRIT}\)&selección de hoja\\
\(F_{\rm obs}\)&\(\mathcal Q_{\rm obs}\to\mathcal Q_{\rm obs}\)&sucesor observable\\
\((\rho_9,q_9)\)&\(\mathbb Z_{>0}\to\mathcal D_9\times\mathbb Z_{\ge0}\)&fase y altura\\
\((r_3,c_3)\)&\(\mathbb Z\to\mathsf{TRIT}\times\mathbb Z\)&orientación y carry\\
\((q_{1000},r_{1000})\)&\(\mathbb Z\to\mathbb Z\times\{0,\ldots,999\}\)&bloque decimal y cociente\\
\(Q\)&\(\mathsf{TPK}\to\mathcal O_{9,12}\)&fase y longitud\\
\(C_{36}^{\rm pred}\)&cociente de \(X_{9,12}\)&memoria predictiva mínima\\
\(C_{108}\)&calendario cíclico&extensión no escindida\\
\(q_{\rm mem}\)&\(\mathbb Z_{\ge0}\)&memoria vertical no acotada\\
\(\operatorname{Ext}_r\)&
\(\operatorname{Ext}_r\subseteq\mathcal F_q\times\mathcal F_{q'}\)&
prolongaciones compatibles\\
\bottomrule
\end{array}
\]

\subsection{Obstrucciones de los lectores y del calendario del TPK: cocientes euclídeos, \(C_{36}^{\mathrm{pred}}\), extensión \(C_{108}\) y memoria vertical}

\begin{criterion}[Criterios de refutación de lectores y calendario]
La unidad queda refutada si:
\begin{enumerate}
 \item un lector residual se presenta como transición completa;
 \item se omite el cociente de alguna división euclídea utilizada después;
 \item se identifica \(C_{36}^{\rm pred}\) con \(R_{36}\);
 \item se afirma que la extensión \(C_{108}\) es el producto
       \(C_{12}\times C_9\);
 \item se declara retorno completo para un número no divisible por \(108\);
 \item el retorno de \(\beta\) fuerza \(q_{\rm mem}=0\);
 \item \(\mathcal K_{27}\) se presenta como un reloj distinto y no como
       una potencia del transporte.
\end{enumerate}
\end{criterion}

\section{Estructura operatorial completa de APP y jerarquía aritmético--geométrica}

El desarrollo operatorial de APP fija la jerarquía
\(4\to2\to6\to54\), los proyectores aritméticos, el bloque central y la
geometría local de caminos. Estas consecuencias permanecen subordinadas al
núcleo formal ya construido y no introducen un segundo origen de APP.

La unidad siguiente construirá la subdivisión cuádruple, el funtor
ciclotómico y los canales internos de noventa y ciento veinte posiciones.
Sólo después se introducirá el espacio exhaustivo de semillas.
